\chapter{图像识别与检索}
\label{chap:vision-image-recognition}

一张路口照片中有两辆车。给出“街道”这个场景名，尚未说明车辆在哪里；给出两个框，尚未说明哪辆是红色、是否存在遮挡、距离相机多远；找到另一张相似街景，也未必找到了同一辆车。本章围绕这些逐渐具体的请求，建立从像素到语义、范围、测量和对应关系的计算。

通用的卷积、注意力和表示学习已经提供了特征提取工具。在这里，需要进一步说明特征对应整幅图、一个对象还是一个局部位置，监督约束哪个输出，以及模型预测如何还原成用户需要的结果。街景贯穿对象与空间任务，正常零件与划痕则用于说明没有完整缺陷类别表时怎样判断异常。以下算例均为教学构造，数值用于复核机制。

\section{语义与属性判定}
\label{sec:vision-recognition-semantics}

类别说明“属于什么”，属性说明“具有什么特征”，异常判断说明“是否偏离规定的正常状态”。三种输出可以来自同一表示，但需要不同的标签和判定依据。场景分类把整幅画面的环境概括为街道或厨房，对象分类则需要明确被判断的对象；厨房这一场景名不能展开为一份完整的物品清单。

\subsection{对象类别与场景判定}
\label{sec:vision-recognition-classification}

在单标签任务中，一个样本只对应类别表中的一项；多标签任务允许多项同时成立。例如车辆裁剪图可按主要车型给一个标签，整张街景却可以同时含公交车、轿车和行人。损失函数必须遵守这一区别：互斥类别可用softmax，多种可共存对象通常分别输出是否存在的得分。

\subsubsection{已有视觉表示与分类头：DINOv3的任务接入}
\label{sec:vision-recognition-dinov3}

已有编码器可以把图像转成全局向量及保留位置的特征。DINOv3提供了这类可复用表示；本节关心的是拿到表示后怎样完成有标签的任务，其完整自监督训练另属表示学习。全局向量压缩整图信息，空间特征仍与图像块位置相连，二者适合接入不同的预测头\scite{vision-2508-10104}。

一条分类训练记录包含图像$\bm{X}_i$及类别$y_i$。先按所选检查点约定缩放、裁剪和归一化，保持标签仍对应被保留内容，再从冻结编码器提取全局向量；局部部位任务应换成部位裁剪或已对齐的空间特征。编码器权重来自已有预训练，新增线性头可以从小随机权重和零偏置开始。训练、调参和测试图像应按车辆、拍摄批次或物种采集条件分开，避免同一对象的近重复图同时提供学习和测试依据。

设冻结编码器$f_{\bm{\theta}}$得到$\bm{z}=f_{\bm{\theta}}(\bm{X})\in\mathbb{R}^{d}$，分类头包含$\bm{W}\in\mathbb{R}^{k\times d}$与$\bm{b}\in\mathbb{R}^{k}$。对$k$个互斥类别，计算
\begin{equation}
 p(y=r\mid\bm{X})=
 \frac{\exp(\bm{w}_r^\top\bm{z}+b_r)}
 {\sum_{j=1}^{k}\exp(\bm{w}_j^\top\bm{z}+b_j)}\text{，}
 \label{eq:vision-linear-head}
\end{equation}
这里$\bm{w}_r^\top$是$\bm{W}$的第$r$行。负对数给真实类别低概率时较大代价；在样本独立的条件模型中，连乘似然取负对数便成为逐样本和，经验目标为$L=-n^{-1}\sum_i\log p(y_i\mid\bm X_i)$，平均使目标不随批次大小直接放大。冻结训练只更新$\bm{W},\bm{b}$；全量适配还更新$\bm{\theta}$，所需梯度存储和训练数据条件随之变化。

对一个样本，记类别得分为$a_r=\bm{w}_r^\top\bm{z}+b_r$、概率为$p_r$，则$\partial(-\log p_y)/\partial a_r=p_r-\mathbf{1}\{r=y\}$。链式求导给出头部梯度$(p_r-\mathbf{1}\{r=y\})\bm{z}$及偏置梯度。每批求平均后，最简单的梯度更新为$\bm{W}^{(t+1)}=\bm{W}^{(t)}-\eta_t\nabla_{\bm W}L$，偏置同样更新，$\eta_t>0$是学习率；带动量的优化器另外保存梯度统计。按事先规定的更新预算训练，依据不参与梯度的验证集选择头部参数，不声称神经网络已经找到全局最优。后文的任务网络沿用“构造监督、前向计算误差、反向更新、验证选择”的关系，但可训练范围、误差项和输出恢复分别说明。

假设三类车型的得分为$(2,1,0)$。经过softmax计算，概率依次约为0.665、0.245和0.090。它们是在给定类别集合内归一化的结果，不能证明图像中一定存在可辨车辆。若编码器读取整幅街景，得分还可能依赖道路背景；若任务要判定左车的车型，应先取得左车区域，或使用明确的区域监督。

若这条样本真实类别为第二类，得分梯度约为$(0.665,-0.755,0.090)$。下降更新提高第二类相对于其他类的得分；这一方向仅说明当前监督如何作用，并不保证下一批或新领域的分类正确。使用时保留编码器、最终分类头、类别顺序和相同预处理，计算得分后取最大项；需要未知状态时再使用验证确定的拒绝规则。场景分类保留整图，逐车分类则保留对象编号及裁剪来源，两者不能混用同一输出记录。

外观接近的物种分类揭示了另一种困难：全局背景足以区分拍摄地点，却可能不足以区分鸟喙或叶缘。iNaturalist包含细粒度类别和不均衡样本，评价时应同时检查小部位分辨率和类别覆盖\scite{vision-1707-06642}
例如常见类别有90张且全对，少见类别有10张且全错，样本平均准确率为90\%，两个类别的平均准确率却为50\%。这两项分别强调观测频率与类别公平覆盖，不能互相替代。

\subsubsection{文字指定类别与零样本识别：CLIP接口}
\label{sec:vision-recognition-clip}

如果使用时才给出“轿车、货车、公交车”，固定输出维度的分类头不能直接容纳新类别。图文表示把类别描述也编码为向量，使图像可以与使用时给定的候选比较。CLIP的双编码器训练已在\ref{sec:vision-models-contrastive}节展开；这里使用已经学到图文对应的两侧编码器及投影，不为当前类别表再训练一个分类头。

令$\bar{\bm{z}}$与$\bar{\bm{t}}_r$分别为单位长度的图像向量和第$r$条描述向量，按$s_r=\bar{\bm{z}}^\top\bar{\bm{t}}_r$排序。使用温度$\tau>0$可计算候选内归一化得分
\begin{equation}
 q_r=\frac{\exp(s_r/\tau)}
 {\sum_{j\in C}\exp(s_j/\tau)}\text{，}
 \label{eq:vision-clip-classification}
\end{equation}
其中$C$是这次输入的非空候选索引集合，图文归一化要求原向量非零。分母随候选变化，$q_r$不是不依赖候选的存在概率。可以对“这是一张公交车照片”等多个模板的单位文字向量求平均后再归一化，但模板集合须在评价前固定，并排除平均后为零的退化情形。

设四项相似度为$0.31,0.28,0.19,0.12$，最后一项是不在图中的“帆船”。最大值只使第一项在这四项中领先。把前三项全部换成不相关词，模型仍会返回一个最大值。开放使用要另设含未知对象的验证协议、拒绝阈值或目标存在判断；单凭相似度最大不能完成拒识。整图相似也无法保证“红色”绑定到货车而非旁边的轿车。

\subsection{对象属性判定}
\label{sec:vision-recognition-attributes}

要记录“红色货车”和“蓝色轿车”，仅输出整图属性集合$\{\text{红色},\text{蓝色}\}$会丢掉归属。可以先定位对象，再对各区域提取表示，形成$(\text{对象编号},\text{属性},\text{取值})$记录。位置与属性共享编号，使后来计数、搜索或人工核验仍指向同一个对象。

属性之间是否互斥应由标签定义决定。“红色车身”和“车灯开启”可以同时成立，“完整车身的主色”则可以规定为一个单标签字段。对于可共存的$k$个二值属性，令$p_{i,r}$为预测，$a_{i,r}\in\{0,1\}$为标注，$m_{i,r}$表示该属性是否有可靠标注，则可用
\begin{equation}
 L_{\mathrm{attr}}=
 -\frac{\sum_{i,r}m_{i,r}
 [a_{i,r}\log p_{i,r}+(1-a_{i,r})\log(1-p_{i,r})]}
 {\sum_{i,r}m_{i,r}}\text{。}
 \label{eq:vision-attribute-loss}
\end{equation}
分母要求至少有一个有效标注；全无标注的批次不计算此项。把“未标注”直接置为“不存在”，会训练出系统性的假负例。

具体接入可令区域向量$\bm{z}_i$经过每属性一个线性得分$h_{i,r}$，再用sigmoid得到$p_{i,r}=1/(1+e^{-h_{i,r}})$。单个有效标注的未平均二元目标对得分的梯度是$p_{i,r}-a_{i,r}$，其中$a_{i,r}$仍指二值标注。冻结区域编码器时只训练属性头，也可按数据条件适配编码器。各属性阈值由独立验证记录选择，输出保留每车的属性分数和可见性；若主色被定义为互斥字段，则另用softmax头，不把二元输出强行当成唯一主色。

裁剪区域过小会切掉车灯，过大又可能纳入旁车颜色。应检查属性所需部位是否可见，并在不可判定时保留未知状态。完整对象编号并不能自行修正错误的视觉证据，但它使属性错绑能够被单独定位。

\subsection{正常状态与异常判定}
\label{sec:vision-recognition-anomaly}

工业零件的异常类型可能不断出现，训练集却主要含正常样本。此时可学习或保存正常外观，再度量新观测与正常外观的偏离。重建误差、分布密度和局部近邻距离是不同的偏离依据，得分都需要结合独立验证集解释。

\subsubsection{正常局部特征记忆与近邻评分：PatchCore}
\label{sec:vision-recognition-patchcore}

PatchCore从已有编码器的中间层提取并局部聚合特征，将正常图像中各位置的向量存入记忆。相邻尺度的特征对齐后保留局部上下文，使用时不需要先列出所有划痕和破损类别\scite{vision-2106-08265}

全部正常片段可能很大，因而从特征集合$\bm{M}$中保留代表子集$\bm{M}_{\mathrm{c}}$。其覆盖目标可写为
\begin{equation}
 \min_{\substack{\bm{M}_{\mathrm{c}}\subseteq\bm{M}\\
 |\bm{M}_{\mathrm{c}}|=m}}
 \max_{\bm{z}\in\bm{M}}\min_{\bm{u}\in\bm{M}_{\mathrm{c}}}
 \lVert\bm{z}-\bm{u}\rVert_2\text{。}
 \label{eq:vision-patchcore-coreset}
\end{equation}
原方法使用近似贪心选择，而非精确穷举全部子集。已选一个起点后，反复加入离当前子集最远的特征，可以保留正常模式中的不同部分；这个过程在建立记忆，不是在用缺陷标签重新训练编码器。

选择时为每个正常特征保存到当前子集的最近距离，新加入$\bm{u}$后更新为原距离与$\lVert\bm{z}-\bm{u}\rVert_2$的较小者，达到预算$m$就停止。原工作用随机投影辅助近似选择，保留的片段仍用于原特征空间评分。最内层最小值寻找一个正常参照，外层最大值检查最难覆盖的正常模式；若只最小化平均距离，少量正常模式可能在记忆中完全缺席。完成后保存编码器、所用中间层及聚合设置、代表记忆和验证阈值。建立记忆的数据必须符合“正常”定义，异常混入记忆会把某些缺陷变成近邻参照。

查询图像位置$p$的基本局部分数为$d_p=\min_{\bm{u}\in\bm{M}_{\mathrm{c}}}\lVert\bm{z}_p-\bm{u}\rVert_2$。将这些分数按空间排列、插值和平滑可形成异常热图。最大局部距离$s_*=\max_p d_p$是整图评分的起点；原方法还根据最异常片段$\bm{z}_*$所对应正常近邻$\bm{u}_*$周围的支持情况重新加权。令$\bm{N}_b(\bm{u}_*)$包含$\bm{u}_*$及其在记忆中的近邻，则
\begin{equation}
 s=\left(1-
 \frac{\exp\lVert\bm{z}_*-\bm{u}_*\rVert_2}
 {\sum_{\bm{u}\in\bm{N}_b(\bm{u}_*)}
  \exp\lVert\bm{z}_*-\bm{u}\rVert_2}\right)s_*\text{。}
 \label{eq:vision-patchcore-score}
\end{equation}
邻域大小$b$是需固定的设置，实际应大于1；否则分子分母相同，分数退化为零。若最近正常片段较孤立，其他邻居的距离更大，括号中的系数更接近1，最大距离受到的压低更少。这个重加权改变整图分数，不改变局部距离图的来源。

假设一维正常记忆为$\{0,1,4\}$，查询局部特征为$\{0.2,0.9,6\}$，对应距离是$0.2,0.1,2$，最大值由第三个片段给出。若$b=2$，取最近正常片段4及其邻居1，整图分数为$2[1-\exp(2)/(\exp(2)+\exp(5))]\approx1.905$。区域线索保留了小范围的偏离，但若新设备使所有特征整体平移，正常零件也可能远离记忆。高分只能表示对当前正常参照的偏离，不能直接命名缺陷或判定业务不合格。

\section{对象定位与区域划分}
\label{sec:vision-recognition-regions}

定位要求把判断接回图上的位置。边界框给出大致范围，掩码给出像素归属，透明或混合边界还需要连续贡献比例。希望处理哪个目标，可以由固定类别、文字或点框提示指定；目标选择与范围形成需要分别解释。

\subsection{对象范围的定位}
\label{sec:vision-recognition-detection}

\term{边界框}（bounding box）用一个矩形包围目标。以下采用连续边界$\bm{b}=(x_1,y_1,x_2,y_2)$，宽高分别为$x_2-x_1$和$y_2-y_1$，不加离散像素计数中的1。检测结果是多个$(\bm{b},c,s)$组成的集合，其中$c$是类别，$s$是得分。标签应说明框包围可见部分还是完整对象；旋转目标也可改用中心、宽高和角度表示。

比较两个框的重叠程度时，\term{交并比}（intersection over union, IoU）将交集面积除以并集面积：
\begin{equation}
 \operatorname{IoU}(\bm{b},\bm{b}')
 =\frac{|\bm{b}\cap\bm{b}'|}
 {|\bm{b}|+|\bm{b}'|-|\bm{b}\cap\bm{b}'|}\text{。}
 \label{eq:vision-iou}
\end{equation}
这里矩形应有正面积，无交集时分子为零。取两个面积均为800的框，边界分别为$(10,10,50,30)$和$(20,10,60,30)$。它们的交集面积为600，IoU为$600/1000=0.6$。若它们是对同一辆车的两个预测，应去重；若属于两辆互相遮挡的车，直接去掉一个则会漏检。

\subsubsection{整图网格预测：YOLO}
\label{sec:vision-recognition-yolo}

原始YOLO把图像划为$S\times S$网格，每个单元预测$B$个框和一组共享的条件类别概率。目标中心落在哪个单元，就由该单元承担监督。框中心相对于单元定位，宽高相对于整图归一化；每框置信度学习目标存在与框质量的联合信号\scite{vision-1506-02640}

若单元列行为$(j,i)$，预测中心偏移为$(u,v)$，输入宽高为$(w,h)$，则中心还原为$((j+u)w/S,(i+v)h/S)$。训练在负责单元的多个预测框中选择当前IoU最高者承担该对象的坐标监督。原始平方误差目标分别约束中心、宽高的平方根、置信度和类别；没有负责对象的候选也接受置信度监督，但权重较低，以免大量背景压倒对象信号。

为看清同一目标怎样进入误差，记单元为$a$、框槽为$b$，$I^{\mathrm{resp}}_{a,b}$标出当前负责框，$I^{\mathrm{other}}_{a,b}$标出其他框槽，$I^{\mathrm{obj}}_a$标出含目标中心的单元，三者均为零一指示。预测中心、宽高和置信度分别为$(\hat x,\hat y)$、$(\hat w,\hat h)$和$\hat C$，真实值去掉帽子；宽高为正并相对整图归一化。$\lambda_{\mathrm c},\lambda_{\mathrm n}>0$分别是坐标和非负责框置信误差的权重。先把位置误差分为中心与尺寸两部分：
\begin{equation}
 \begin{aligned}
 e_{a,b}^{\mathrm{pos}}&=(\hat x_{a,b}-x_a)^2
                       +(\hat y_{a,b}-y_a)^2,\\
 e_{a,b}^{\mathrm{size}}&=(\sqrt{\hat w_{a,b}}-\sqrt{w_a})^2
                        +(\sqrt{\hat h_{a,b}}-\sqrt{h_a})^2,\\
 L_{\mathrm{coord}}&=\lambda_{\mathrm c}\sum_{a,b}
 I^{\mathrm{resp}}_{a,b}(e_{a,b}^{\mathrm{pos}}+e_{a,b}^{\mathrm{size}})\text{。}
 \end{aligned}
 \label{eq:vision-yolo-training-process}
\end{equation}
类别和置信误差则按同一职责选择监督：
\[
 \begin{aligned}
 L_{\mathrm{conf}}&=\sum_{a,b}I^{\mathrm{resp}}_{a,b}
 (\hat C_{a,b}-\operatorname{IoU}_{a,b})^2\\
 &\quad+\lambda_{\mathrm n}\sum_{a,b}
 I^{\mathrm{other}}_{a,b}\hat C_{a,b}^{\,2},\\
 L_{\mathrm{class}}&=\sum_a I^{\mathrm{obj}}_a\sum_r
 (\hat p_a(r)-\mathbf{1}\{r=c_a\})^2\text{。}
 \end{aligned}
\]
平方根使相同绝对宽高误差在较小框上占更大相对权重；它不改变中心仍由单元承担的分配。对训练图求平均，将三项相加，依当前预测重新选择负责框，再反向更新网络。原工作先做分类预训练再训练检测网络，新增检测层接受框监督；推理只保留最后的网络，不保留逐批负责关系。

教学中$S=7$、输入$448\times448$，单元列行$(2,3)$的中心偏移$(0.5,0.25)$还原为$(160,208)$。若归一化宽高为$(0.2,0.1)$，像素宽高是$(89.6,44.8)$，再由中心加减半宽高得到四边。置信度0.8与该单元货车条件得分0.75相乘，货车框分数为0.6。输入曾缩放或加边时，最后还须去掉加边并除缩放比例，才是原图候选。

每框置信度写为$\mathbb{P}(\mathrm{obj})\operatorname{IoU}$，乘上共享的$\mathbb{P}(c\mid\mathrm{obj})$得到类别相关得分。它将语义与定位质量合在一起，不能只解释为类别概率。使用时解码全部框，按阈值过滤，再用\term{非极大值抑制}（non-maximum suppression, NMS）去掉同类中与当前高分框过度重叠的低分框。

NMS按得分从高到低取一个框，保留它并删除重叠超过阈值的候选，再对余下框重复，直到没有候选。对于前述IoU为0.6的两个框，若得分为0.9与0.7，阈值为0.5，第二个会被删除。这一规则处理重复预测，却可能伤害拥挤场景。一个单元共享类别预测也限制了原始版本处理同单元多对象的能力；后来的YOLO版本改变了预测、分配与去重设计，不能沿用此处全部细节。

\subsubsection{多尺度特征与密集训练权重：FPN、RetinaNet}
\label{sec:vision-recognition-fpn}

远处车辆在图上只占几个像素，深层特征的空间步幅可能使它消失。浅层保留更多位置细节，深层又具有较大的上下文。\term{特征金字塔网络}（feature pyramid network, FPN）把深层特征上采样，与相应浅层经通道变换后的特征相加，再用卷积整理，形成不同空间尺度的预测特征\scite{vision-1612-03144}

第$\ell$层的步幅$s_\ell$连接特征坐标与原图坐标。训练需要按使用该金字塔的检测器规定，将目标或锚框分配到尺度层；FPN本身只生成特征，并不规定一种通用的最终检测集合。较密特征有助于保留小目标，但若目标在输入缩小时已经丢失，特征融合不能补回观测。

记主干第$\ell$层为$\bm{C}_\ell$，自顶向下的中间特征为$\bm{M}_\ell$，则相同空间网格上可以写$\bm{M}_\ell=\operatorname{Conv}_{1\times1}(\bm{C}_\ell)+\operatorname{Up}(\bm{M}_{\ell+1})$，再用$3\times3$卷积得到$\bm{P}_\ell$。横向卷积先统一通道数，上采样先统一位置数，否则两项无法相加。对步幅$s_\ell$的输出位置$(i,j)$，检测器按其网格约定放置原图参考中心；不同层即使特征形状不同，最终框也需回到同一原图坐标再比较。

RetinaNet在金字塔各位置放置预定大小和长宽比的\term{锚框}（anchor box），即坐标回归的参考框。网络预测类别以及相对于参考框的中心偏移、宽高对数修正。训练按锚框与真实框的IoU分配前景、背景和忽略项，推理则解码修正、筛选并执行NMS。

设锚框中心和尺寸为$(x_a,y_a,w_a,h_a)$，目标为$(x_*,y_*,w_*,h_*)$。训练回归标签为$((x_*-x_a)/w_a,(y_*-y_a)/h_a,\log(w_*/w_a),\log(h_*/h_a))$；给定预测修正$\bm{\delta}$，逆变换为
\begin{equation}
 \begin{aligned}
 \hat x&=x_a+w_a\delta_x,\qquad \hat y=y_a+h_a\delta_y,\\
 \hat w&=w_a e^{\delta_w},\qquad \hat h=h_a e^{\delta_h}\text{。}
 \end{aligned}
 \label{eq:vision-retinanet-decode}
\end{equation}
中心差除参考尺寸，宽高比取对数，使标签无像素单位且能表达乘性缩放。教学锚框$(50,30,40,20)$配目标$(60,30,80,20)$，标签为$(0.25,0,\log2,0)$，逆算恢复原目标。负锚框没有对应框，不计算坐标回归项；边界分配的忽略项不参与分类。

大量背景锚框在训练后很容易正确分类，但数量仍可能主导总损失。令$p_t$为模型分给真实二值标签的概率，\term{焦点损失}（focal loss）为
\begin{equation}
 L_{\mathrm{focal}}=-\alpha_t(1-p_t)^\gamma\log p_t\text{，}
 \qquad\gamma\geq0\text{。}
 \label{eq:vision-focal-loss}
\end{equation}
其中$\alpha_t$调整正负样本权重。$\gamma=0$时退回加权交叉熵；$\gamma=2$时，$p_t=0.99$的容易样本乘上$10^{-4}$，而$p_t=0.5$的困难样本乘上$0.25$。这个差别让训练更关注尚未学好的候选\scite{vision-1708-02002}
改变容易背景的权重没有增加少见类别的观测，因而前景背景失衡与类别长尾仍须分别检查。

各锚框按类输出独立sigmoid得分，分类采用焦点项，正锚框回归采用平滑$L_1$项，训练按正例数量归一化；没有正例的批次应按实现保留非零归一化下界。主干可从分类预训练初始化，金字塔与任务头由检测监督学习。对同一批的前景和背景计算误差后更新共享参数，使用时从所有层解码参考框、合并同类候选再NMS，最后逆变换到原图。焦点因子只是改变每个候选的训练贡献，没有从后处理集合中直接删除背景。

\subsubsection{逐位置边界距离预测：FCOS}
\label{sec:vision-recognition-fcos}

预定义参考框带来了大小与比例的设计。FCOS改为让特征位置直接预测到对象左、上、右、下边界的距离$(l,t,r,b)$。若位置对应原图$(x,y)$，预测框就是$(x-l,y-t,x+r,y+b)$，不再以锚框为坐标起点\scite{vision-1904-01355}

监督仍然需要分配。位于真实框内且满足尺度范围的位置可以成为正例；原方法用不同层的回归范围分摊对象，并在多个框竞争一个位置时选择面积较小的框。负例只学习分类，正例还学习距离及中心度。这里描述的是原始FCOS，后续采用中心区域采样等改动时应另记配置。

每层特征位置先按步幅对应原图采样中心，再由目标四边减该中心形成$(l,t,r,b)$标签。训练分类采用焦点目标，正例的恢复框与真实框用IoU误差比较，中心度使用二元交叉熵；三项按权重组成批次目标，梯度同时到达特征主干及三个预测分支。距离分支按实现约束为正，推理由距离和采样中心直接得到四边，类别分数与中心度组合后跨层筛选。对于无目标位置，只学“不是这些类别”，不能给它虚造四边距离。

\term{中心度}（centerness）衡量位置是否接近框中心，训练目标为
\begin{equation}
 c_*=\sqrt{\frac{\min(l,r)}{\max(l,r)}
                 \frac{\min(t,b)}{\max(t,b)}}\text{。}
 \label{eq:vision-fcos-centerness}
\end{equation}
式中距离来自分配的真实框，框内部距离为正。中心点四侧相对均衡，$c_*=1$；靠近边界时其中一个比例趋于零。推理把预测中心度与类别得分组合，降低边缘候选的排序，再用NMS形成集合。

对于框$(20,10,80,50)$，中心$(50,30)$的目标距离为$(30,20,30,20)$。位置$(25,30)$的距离变成$(5,20,55,20)$，中心度为$\sqrt{1/11}\approx0.302$。即使这两个位置都能精确回归同一框，中心度也不同，说明它不是预测框IoU的直接测量。图\ref{fig:vision-box-points}把位置、距离与框的关系画在同一坐标系中。

\begin{figure}[htbp]
 \centering
 % 连续像素坐标下的FCOS教学例；所有长度与正文算例一致。
\begin{tikzpicture}[x=1mm,y=-1mm,font=\figurefont\small,
  dist/.style={-{Stealth[length=2mm]},draw=BookBlue,line width=0.9pt}]
 \draw[draw=BookMuted,line width=0.7pt,->] (10,0)--(90,0) node[right] {$x$};
 \draw[draw=BookMuted,line width=0.7pt,->] (10,0)--(10,58) node[below] {$y$};
 \draw[draw=BookTeal,fill=BookTeal!7,line width=1pt] (20,10) rectangle (80,50);
 \draw[draw=BookMuted,dashed,line width=0.7pt] (50,10)--(50,50);
 \draw[draw=BookMuted,dashed,line width=0.7pt] (20,30)--(80,30);
 \fill[BookInk] (50,30) circle (1mm);
 \fill[BookBlue] (25,30) circle (1mm);
 \node[anchor=north west] at (51,33) {中心：$c_*=1$};
 \node[anchor=south west] at (26,27) {边缘：$c_*\approx0.302$};
 \draw[dist] (25,18)--(20,18);
 \draw[dist] (25,18)--(80,18);
 \node[anchor=south] at (22.5,17) {$5$};
 \node[anchor=south] at (57,17) {$55$};
 \draw[draw=BookBlue,line width=0.7pt] (25,18)--(25,29);
 \node[anchor=north] at (20,51) {$(20,50)$};
 \node[anchor=south] at (80,9) {$(80,10)$};
\end{tikzpicture}

 \caption{同一真实框内的中心位置与边缘位置。实线箭头标出边缘位置到左右边界的距离，虚线标出中心轴；两点都可以回归同一矩形，但中心度目标分别为1与约0.302。坐标采用向右、向下为正的连续像素边界。}
 \label{fig:vision-box-points}
\end{figure}

\subsubsection{无序对象集合与唯一匹配：DETR}
\label{sec:vision-recognition-detr}

密集候选需要在输出端去重。DETR把目标直接视为无序集合：固定$n$个对象查询读取图像特征，每个查询预测一个框和类别，类别还包含“无对象”。查询是学习得到的向量，并非预先指定给某辆车的身份\scite{vision-2005-12872}

训练标注有$m\leq n$个对象，记第$i$个真实类别与框为$(c_i,\bm{b}_i)$，第$j$个预测为$(\hat{\bm{p}}_j,\hat{\bm{b}}_j)$。寻找从真实对象到不同预测位置的单射$\sigma$：
\begin{equation}
 \hat{\sigma}=\argmin_\sigma
 \sum_{i=1}^{m}\left[
 -\hat p_{\sigma(i)}(c_i)
 +\lambda_1\lVert\bm{b}_i-\hat{\bm{b}}_{\sigma(i)}\rVert_1
 +\lambda_{\mathrm{g}}L_{\mathrm{GIoU}}
    (\bm{b}_i,\hat{\bm{b}}_{\sigma(i)})\right]\text{。}
 \label{eq:vision-detr-matching}
\end{equation}
框在共同归一化坐标中比较。广义IoU还惩罚两个框最小包围矩形中未被并集覆盖的区域，使无交集时也有几何信号。匹配代价用于选对应；选定后，分类损失采用负对数概率，定位项只计算匹配对象，未匹配查询学习“无对象”，并降低该类的权重。

具体计算中，主干从缩放图像产生二维特征，展平后加位置编码，编码器在图像位置间交换上下文。每个可学习查询经解码器读取位置特征，再由共享类别头和框头产生结果。框头输出归一化中心、宽高；预测序列没有必须对应目标编号的顺序。匹配可以用匈牙利算法求当前代价矩阵的最小单射，选定后视对应为本轮固定监督，不对离散匹配本身求梯度。

若$E$为两框最小包围矩形，广义交并比为$\operatorname{GIoU}=\operatorname{IoU}-|E\setminus(\bm b\cup\bm b')|/|E|$，定位损失使用$1-\operatorname{GIoU}$。未相交框的IoU均为零，这个额外空白比例仍能区分隔得近和隔得远。训练目标是匹配查询的负对数分类及两项框误差，加未匹配查询的无对象分类；多图求平均后更新主干、编码解码器、查询和任务头。原始实现可用已有分类主干初始化，其余部分由检测任务学习。

假设两个真实对象对三个查询的总代价矩阵为
\[
 \bm{C}=\begin{pmatrix}1&2&9\\1.2&5&8\end{pmatrix}\text{。}
\]
逐行独立取最小都会选第一个查询，违反唯一对应。集合匹配选择第一对象到第二查询、第二对象到第一查询，总代价为$3.2$，小于反过来分配的$6$。三辆车的情形同样要求三个不同查询；改变标注排列只改变匹配的表示，不应改变最佳集合代价。

使用时没有真实框参与匹配，直接过滤“无对象”和低分输出。原始DETR不执行常规NMS，重复预测主要由唯一匹配训练约束。查询数限制一次可输出的对象数，原始全局注意力、训练周期与小目标效果也各有条件；这些限制不能概括后来所有Transformer检测器。

\Needspace{65mm}
\subsubsection{文字条件下的对象预测：Grounding DINO}
\label{sec:vision-recognition-grounding-dino}

当目标变成“红色货车”时，需要把类别与属性文字接入检测计算。Grounding DINO分别编码图像和文字，通过跨模态特征增强交换信息，再用语言相关性选择用于初始化的对象查询；解码器继续读取两侧特征并更新框与文字对应得分\scite{vision-2303-05499}

文字经分词形成若干词元，候选图像特征与词元比较，训练用对象范围以及对象对应的词语监督匹配。推理时不仅要还原框，还要说明得分对应“货车”、包含属性的短语，还是更长表达中的一部分。词元阈值、框阈值和短语聚合规则都会影响最终输出。

令增强后的候选图像特征为$\bm{v}_i$、文字词元为$\bm{t}_r$，语言引导查询选择依据候选与文字的相似度，例如从$\max_r\bm{v}_i^\top\bm{t}_r$较高的位置选初始化查询。此时选中的仍是候选位置，不能当成已经确认的货车。跨模态解码继续更新查询和归一化框，并形成各查询对各词元的匹配logit。训练记录需包含图像、文本、目标框及文本中的目标词元关联，关联可以跨多个词元；分类误差约束这些关联，$L_1$和广义IoU约束匹配框，集合分配避免同一监督对象任意占多个查询。

可使用已有视觉和语言编码器初始化，经过有图文范围对应的预训练或任务适配获得连接权重；使用已发布检查点时保留其分词、图像尺度及阈值协议，不为每次文字查询更新权重。以“红色货车”查两辆车时，先解码每个候选框，再保留满足框阈值的项，从相应词元得分恢复短语，最后将框中心宽高乘原图尺寸并转换为四边。若模型只在“货车”词元高分而“红色”得分弱，最终记录不能悄悄把整个属性短语写成已核实事实。

“红色货车”要求属性绑定，“公交车旁的货车”又要求利用参照物消歧。能接收这两种字符串只说明接口容纳了条件，不能证明关系已经正确解析。可分别遮去参照物、更换颜色词或保留多个同类目标，检查预测是否按所需证据改变。完整的表达定位及证据协议在\ref{chap:vision-multimodal-understanding}章讨论。

\subsubsection{开放词汇定位之外的未知对象发现}
\label{sec:vision-recognition-unknown-discovery}

仓库画面中出现一个尚未加入类别表的搬运装置。固定分类器可能把它认作货车，也可能把它当背景；文字定位允许用户写出装置名称，但用户此时未必知道应当查询什么。\emph{未知对象发现}（unknown object discovery）要求在没有名字时先提出“这里可能有一个对象”，保留区域供后续核验。它与开放词汇定位的差别在于名称是否已作为输入，也与拒识不同：拒识已有裁剪图只回答是否属于已知类，发现还要确定未知对象在哪里。

用集合$C_{\mathrm{k}}$表示训练已标注类别，输出可以是已知记录$(\bm{b},c,s)$，也可以是未命名记录$(\bm{b},\mathrm{unknown},s_{\mathrm{o}})$。对象性（objectness）$s_{\mathrm{o}}$衡量一个候选是否应被当作独立对象，而不是它属于哪个已知类。未匹配标注的候选中同时含真实背景和未标注对象，把它们全部用零监督，会使发现能力向已知类别偏移。

OWOBJ把对象性建模接入已有检测器的训练，而不另造一套推理接口。以其基于可变形DETR的实现为例，图像特征经编码、查询解码得到$\bm{q}_i\in\mathbb{R}^d$；已知类别头、框头和对象性分支读取这些查询。训练先按已知框和类别作唯一匹配，匹配查询确定属于对象，其他查询暂不能判成背景。原论文由查询特征的移动平均估计均值$\bm{\mu}$与正定对角协方差$\bm{\Sigma}$，用
\begin{equation}
 \begin{aligned}
 d_{\mathrm{M}}(\bm{q})=
 (\bm{q}-\bm{\mu})^\top\bm{\Sigma}^{-1}(\bm{q}-\bm{\mu}),\\
 s_{\mathrm{o}}(\bm{q})=\exp[-d_{\mathrm{M}}(\bm{q})]\text{。}
 \end{aligned}
 \label{eq:vision-owobj-objectness}
\end{equation}
生成对象性软标签。马氏距离（Mahalanobis distance）按各特征方向的正常变动尺度比较偏离，避免原始数值尺度大的通道独占距离。$s_{\mathrm{o}}$在距离为零时为1，距离增大时下降；它是该训练设计中的对象性信号，不是已经校准的真实对象概率\scite{vision-owobj-2025}。

匹配查询继续接受已知类别和框监督，并通过减小$d_{\mathrm M}$提高对象特征的密度；未匹配查询用$s_{\mathrm{o}}$而非一律为零的标签监督额外对象性输出。原工作还用变分近似的先验匹配项约束分布：高斯先验的方差随估计后验调整，并加入方差偏置，以处理少标注条件下的分布收缩。已知与未匹配查询之间的能量间隔项则改变分类得分的相对约束。这里的能量来自分类logit的对数指数和，不是像素亮度或物理能量。训练更新对象性相关特征及检测器参数，移动统计随批次更新；使用阶段沿用所接检测器的框和对象性输出，不再用真实标注给新图匹配。

例如批次的经验均值为$\bar{\bm q}$，移动均值可按$\bm\mu^{(t+1)}=\rho\bm\mu^{(t)}+(1-\rho)\bar{\bm q}$更新，$0\leq\rho<1$控制旧批次保留程度，对角方差同样利用当前统计更新。方差需要保持正数才能求逆，初始化和数值下界必须随实现固定。整体训练在已有检测框与分类误差外加入匹配对象的距离项、先验匹配及能量间隔，仍从所接检测器的预训练权重和规定预算开始。未知输出筛选沿用该检测器的对象性分类器，而式\eqref{eq:vision-owobj-objectness}主要提供训练软标签；不能在使用时拿新图真标签重新估计对象分布。

作为一维教学例子，若$\mu=0$、$\Sigma=1$，两个未匹配查询为0.2和2，则软对象性分别为$e^{-0.04}\approx0.961$和$e^{-4}\approx0.018$。给前者较强对象监督可以避免直接把未知候选抹掉，但若背景也靠近这个分布，它同样会被提出。这说明对象性需要同时检查未知召回和背景误检，不能由“保留更多区域”推断发现已经正确。

发现后可暂存原图编号、框、得分和当前类别表，再由人工或后续证据命名新类，形成新增训练记录。加入新类别是类别增量（class-incremental learning）条件，即后续学习扩展输出集合还要保留旧类能力；它不是这一次未命名检测已经完成的步骤。增量更新可参考\BookChapterRef{chap:ket-update-maintenance}{卷八“如何更新并维持已有知识”章}，当前必须保留未知区域及其来源。仓库装置被误报为已知货车、墙面被误报未知对象、加入装置类后货车退化，是三个不同错误。概念提示可在\ref{sec:vision-recognition-sam3}节进一步产生实例掩码，但只查询已有名称仍不能替代未知发现协议。

\subsection{像素归属与对象区分}
\label{sec:vision-recognition-segmentation}

框内常含背景，也可能含另一辆车。需要准确范围时，应进一步指定像素归属。\term{语义分割}（semantic segmentation）给像素分配类别，同类车辆可以连成一片；\term{实例分割}（instance segmentation）给不同车辆各自的掩码；\term{全景分割}（panoptic segmentation）还共同覆盖道路、天空等非实例区域，使每个像素获得类别，可数对象的像素再获得实例编号。

这一区别决定结果的组织方式。两辆车的并集可以是正确的车辆语义区域，却是错误的单个实例。道路通常按区域类别处理，不需要把每块路面任意编号。检测、语义与实例标签不能未经说明相互转换。

图\ref{fig:vision-segmentation-outputs}沿用同一个离散街景，显示增加实例编号与背景类别各带来什么信息。这里仅画出可见区域，框状图元表示教学掩码，不是算法检测框。

\begin{figure}[htbp]
 \centering
 % 同一离散场景的三种分割输出；字符在灰度下保留类别和身份区别。
\begin{tikzpicture}[font=\figurefont\small,x=4.5mm,y=-4.5mm]
 \foreach \offset/\title in {0/语义分割,8/实例分割,16/全景分割}{
  \begin{scope}[xshift=\offset*4.5mm]
   \node[anchor=south] at (3,0) {\title};
   \fill[BookPaper] (0,0.4) rectangle (6,4.4);
   \draw[BookMuted,line width=0.7pt] (0,0.4) rectangle (6,4.4);
   \draw[BookRule,line width=0.5pt] (0,1.4)--(6,1.4);
   \ifnum\offset=8
    \node[text=BookMuted] at (3,3.6) {不标背景};
   \else
    \node[text=BookMuted] at (3,3.6) {道路};
    \node[text=BookMuted] at (3,0.9) {天空};
   \fi
   \filldraw[fill=BookTeal!18,draw=BookTeal,line width=0.8pt]
      (0.5,1.8) rectangle (2.8,2.8);
   \ifnum\offset=0
    \filldraw[fill=BookTeal!18,draw=BookTeal,line width=0.8pt]
      (3.2,1.8) rectangle (5.5,2.8);
    \node at (1.65,2.3) {车};
    \node at (4.35,2.3) {车};
   \else
    \filldraw[fill=BookAmber!16,draw=BookAmber,dashed,line width=0.8pt]
      (3.2,1.8) rectangle (5.5,2.8);
    \node at (1.65,2.3) {车1};
    \node at (4.35,2.3) {车2};
   \fi
  \end{scope}
 }
\end{tikzpicture}

 \caption{三种分割结果的区别。语义分割保留类别，实例分割分开可数对象，全景分割共同保留背景类别与对象编号；文字编号和边界线型使身份区别在灰度打印中仍可辨认。图中区域为离散教学构造。}
 \label{fig:vision-segmentation-outputs}
\end{figure}

\subsubsection{从空间特征形成像素判断}
\label{sec:vision-recognition-dense-pixels}

密集特征仍比原图稀疏。把它们简单放大只能增加采样位置，细边界需要由保留位置的信息支持。单幅语义图和双时相变化图都以像素为输出，但前者判断类别，后者判断同一区域的状态是否改变。

\Needspace{60mm}
\paragraph{多尺度上下文与边界恢复：DeepLabv3+}
\label{sec:vision-recognition-deeplab}

若减小空间步幅，计算成本随特征尺寸上升；若持续下采样，路杆和车窗边缘可能难以恢复。DeepLabv3+用空洞卷积扩大特征读取范围，结合不同采样间隔的上下文，再以解码器融合较浅层细节\scite{vision-1802-02611}

空洞率$r$表示卷积核相邻采样位置间隔为$r$，例如一维三点核读取$x_{i-r},x_i,x_{i+r}$，参数个数仍为三。多尺度分支与全局汇总使位置可以同时利用邻域和较大范围。解码阶段上采样语义特征，对浅层特征先调整通道，再拼接和卷积，最后恢复输入网格上的类别得分。

训练对有标注像素计算交叉熵，忽略无效标签；推理逐位置选择类别并按预处理的逆坐标变换恢复原图。道路内部的类别主要依赖上下文，细路杆又需要足够分辨率。改变输出步幅和输入尺寸会改变这两类证据与预算，应共同记录。即使两车边缘都清楚，纯类别图仍未分配实例身份。

对训练图的有效位置集$\Omega_i$，类别标注为$y_{i,p}$，输出logit为$h_{i,p,r}$。按类别softmax得到$p_{i,p,r}$后，像素目标为
\begin{equation}
 L_{\mathrm{pixel}}=
 -\frac{\sum_i\sum_{p\in\Omega_i}\log p_{i,p,y_{i,p}}}
        {\sum_i|\Omega_i|}\text{。}
 \label{eq:vision-dense-pixel-training}
\end{equation}
每个位置的得分梯度是预测类别分布减真实类别指示，再经上采样、解码器和上下文分支传回主干。可从分类主干初始化，新增上下文和解码层由像素监督学习。缩放、裁剪图像时，对离散标签采用保持类别的采样，不能用颜色插值制造新类别。推理宜先把类别得分图插值回原图，再作类别选择；直接插值整数类别编号没有概率含义。该恢复只补采样网格，不补被输入删掉的细杆。

\paragraph{双时相上下文与变化图：BIT}
\label{sec:vision-recognition-bit}

遥感影像的目标可能是发现新增建筑，而不是给单幅影像中的每栋楼分类。BIT输入同一区域的两期图像，以共享的特征提取器形成空间表示，再把每期特征通过学习的空间权重汇总为少量语义词元。两期词元共同编码，使一幅图的局部判断能参考另一时期的上下文\scite{vision-2103-00208}

可将词元汇总理解为$\bm{T}=\bm{A}^{\top}\bm{F}$：$\bm{F}\in\mathbb{R}^{n\times d}$是$n$个位置的特征，$\bm{A}\in\mathbb{R}^{n\times k}$的每列在位置上归一化，得到$k$个加权汇总。编码后的词元再通过交叉注意力返回两期空间位置，恢复上下文增强的特征；对两期特征作差异计算并用预测头输出变化概率。

训练使用配准影像与像素变化标签，模型参数由变化损失更新。新建筑应对应语义变化，树木变绿是否计为变化则取决于标注定义。若两期建筑整体错开几个像素，差异分支也会产生边缘响应，因此配准是任务前提。上下文交换能够利用关联，却不能把传感器差异自动解释成真实建设活动。

每次训练同时裁剪、翻转两期图像及变化标签，保持同一位置仍指同一地表；不能对两期各取互不对应的随机区域。若增强特征分别为$\bm{F}_1',\bm{F}_2'$，变化头从$|\bm{F}_1'-\bm{F}_2'|$形成两类得分，再在有效像素上用式\eqref{eq:vision-dense-pixel-training}的类别目标学习“变化／不变”。主干可以从已有图像权重初始化，词元汇总、双期编码和空间解码由变化数据学习。冻结或适配主干需按实际设置记录。使用时保留相同双期顺序和配准，把分数恢复到共同原网格，阈值选出变化区域；区域只说明标签所定义的状态差异，没有自动解释变化原因。

\subsubsection{从对象或区域形成掩码}
\label{sec:vision-recognition-mask-sets}

为了分开同类对象，可以先取得候选对象，再在对象内预测掩码；也可以直接预测一组掩码并匹配实例。用户点框提示又提供了另一种目标选择依据。以下沿这三种接口展开。

\paragraph{区域对齐与实例掩码：Mask R-CNN}
\label{sec:vision-recognition-mask-rcnn}

Mask R-CNN在候选框上并行预测类别、框修正和实例掩码。为了让框内位置与原特征准确对应，\term{区域对齐}（RoIAlign）保留连续坐标，以双线性插值在各小格中采样，再汇总成固定尺寸特征，避免先把框边界粗略取整造成的位移\scite{vision-1703-06870}

候选区域来自区域建议网络（region proposal network, RPN），即在特征位置的参考框上先预测对象性和框修正，形成可能含对象的矩形。RPN的前景背景分配依据真实实例框；候选解码和去重后，第二阶段再按候选与真实框的重叠分配类别、框修正和实例掩码。训练可以用已有分类主干初始化，并共同优化RPN目标和区域分支目标。一个正候选必须记住匹配实例编号，否则即使“车辆”类别正确，仍可能使用另一辆车的掩码标签。

设四个相邻特征值为$f_{0,0},f_{0,1},f_{1,0},f_{1,1}$，连续位置在水平方向偏移$a$、竖直方向偏移$b$，二者均在$[0,1]$内，则插值为
\begin{equation}
 f(a,b)=(1-a)(1-b)f_{0,0}
 +a(1-b)f_{0,1}+(1-a)bf_{1,0}+abf_{1,1}\text{。}
 \label{eq:vision-roialign}
\end{equation}
用数值$(0,4,2,6)$与$(a,b)=(0.25,0.5)$可得$f=2$，既保留横向的四分之一偏移，也保留纵向中点。若直接取最近格点，这些偏移信息就会消失。

原方法对每个类别预测一个二值掩码通道。训练时只对该正候选匹配到的真实类别通道计算像素二元损失，其他类别通道不在此样本上竞争像素归属。使用时取预测类别的掩码，将区域坐标缩放回原框并放回原图。两个重叠车辆由两个候选各自形成掩码，重叠区域的实例归属与类别判断是不同问题。

区域掩码网格的每个目标值由匹配实例在该区域中的前景归属取得，预测经sigmoid后用二元交叉熵约束。总目标将RPN对象性及框误差、区域类别及框误差、正候选掩码误差相加，再沿共享特征反向更新；负候选没有实例掩码目标。推理先按类别与框分数筛选、去重，再对保留实例取相应掩码通道。若一个$28\times28$掩码对应原框$(20,10,80,50)$，归一化区域位置$(\xi,\eta)$回到$(20+60\xi,10+40\eta)$，放置时还需处理框越界和原图尺寸。恢复后的掩码、类别和框共享实例编号，不能分别排序后重新拼接。

区域对齐解决采样量化误差，不能补回候选阶段漏掉的车辆，也不能修正把候选匹配到错误实例的监督。其关键点扩展改变输出目标为有部位语义的位置，不能由掩码准确直接推断关节准确。

\paragraph{掩码查询与局部注意力：Mask2Former}
\label{sec:vision-recognition-mask2former}

先有框再分割并非唯一方式。Mask2Former从多尺度像素特征与查询直接产生掩码集合。每个查询预测类别以及掩码向量，掩码向量与各位置像素嵌入比较，形成该查询的前景概率；训练用集合匹配连接预测与目标，再计算分类、二元掩码和Dice损失\scite{vision-2112-01527}

Dice项度量软预测与真实前景的重叠。对于预测$p_j$和真实标签$m_j$，一种带平滑量$\epsilon>0$的形式为
\begin{equation}
 L_{\mathrm{Dice}}=
 1-\frac{2\sum_j p_jm_j+\epsilon}
 {\sum_j p_j+\sum_j m_j+\epsilon}\text{。}
 \label{eq:vision-dice}
\end{equation}
它与逐像素损失提供不同信号：一个关注区域重叠，另一个约束各位置的预测。训练中的采样点、归一化及匹配权重还需按具体实现固定。

掩码注意力把上一轮预测的前景区域变为后续跨注意力的可读范围。对于查询$i$与位置$j$，在注意力得分中加入屏蔽项$M_{i,j}$：允许读取为0，不允许为$-\infty$，再作softmax。这样每个查询重点读取自己可能对应的区域；多尺度特征与反复更新使区域可以细化。

设第$i$个查询解码成掩码向量$\bm{e}_i\in\mathbb{R}^d$，位置$p$的像素特征为$\bm{f}_p\in\mathbb{R}^d$，掩码概率为$\sigma(\bm{e}_i^\top\bm{f}_p)$。每次查询更新后重新预测掩码并形成下一层读取范围，训练匹配则比较类别和采样掩码误差，不用检测框替代掩码监督。匹配查询承担分类、二元掩码与Dice项，未匹配查询学无对象；训练数据的裁剪同时作用到目标区域。若屏蔽使所有位置都不可读，注意力的分母将失效，实现需恢复可读位置，而不能直接对全$-\infty$求softmax。

语义输出可将各查询的类别概率与该查询掩码概率相乘，对查询求和，再逐像素选类；实例输出保留查询编号，按类别及掩码质量筛选；全景输出还需处理区域竞争、无对象和不可数类别。三种规则使用相同输入街景，却保存不同信息。例如两个车辆查询可以在语义图中合并，实例图中不能合并。最终将分数图恢复原图后再组织区域，避免把裁剪坐标当作原像素。

原工作按任务与数据分别训练，不能把统一架构理解为一份参数无条件通用。错误的早期掩码也可能限制后续读取，因此要检查小对象和被遮挡对象是否在迭代中持续被漏掉。

\paragraph{图像与提示编码、歧义掩码：SAM}
\label{sec:vision-recognition-sam}

有时用户只希望选中图上的一个区域，甚至尚未给它类别名。SAM把图像编码、提示编码与掩码解码分开，同一图像的昂贵编码可以复用，点或框变化时更新提示并运行较轻的解码器。点提示包含位置与前景、背景标签，框提示提供空间范围；已有掩码也可作为稠密提示\scite{vision-2304-02643}

解码器让提示与图像表示交换信息，输出若干候选掩码及质量估计。训练提供图像、目标掩码与采样提示，歧义输出通过多个候选容纳一个提示的不同合理解释。质量预测估计候选与监督目标的重叠程度，服务选择，并非在使用时知道了真实掩码。

图像编码器从已有视觉权重开始，提示位置先归一化到图像坐标，并附正负点或框角类型的学习嵌入；稠密掩码提示经下采样与图像特征对齐。解码器让提示读取图像，也让图像读取提示，输出掩码向量；向量与上采样图像特征逐位置比较，得到低分辨率掩码logit。多个输出中对监督掩码误差较小的一项承担歧义训练的分割目标，原方法结合焦点及Dice项，并用真实候选IoU监督质量头。训练模拟后续点击，在错误区域补提示，参数由提示条件下的残差更新；不能把真实目标直接放进使用时的质量分支。

使用同一街景时先缓存图像编码，点或框改变只更新提示计算及解码。候选logit上采样到输入图，去除加边后恢复原尺寸，再按阈值变成掩码；质量得分可推荐候选，用户意图仍决定选择车门还是整车。交互停止可以由用户确认或预设点击预算控制，不存在仅凭质量头保证正确意图的终止规则。

落在车门上的正点可能指车门或整车。用户选择接近意图的候选，再在漏掉的车尾补正点，或在被误选的路面补负点，才逐步消除歧义。应报告提示来源、点击次数与纠错策略；用真实标注框作提示与由检测器自动给框，具有不同的信息与错误条件。SAM的区域输出本身不自动附带“货车”名称。

\paragraph{概念存在与实例范围：SAM 3}
\label{sec:vision-recognition-sam3}

点提示通常选择一个局部区域，“找出所有校车”则需要判定概念是否出现，并枚举符合概念的实例。SAM 3接收概念短语和带正负标签的图像区域示例，区域视觉特征参与概念表示；随后由对象查询预测框、掩码和局部匹配得分。原始报告将这一接口称为可提示概念分割（promptable concept segmentation, PCS），即用概念而非单个局部点击指定全部符合条件的实例\scite{vision-2511-16719}。

概念是否在整幅图像中出现，与每个候选是否匹配，承担不同职责。共享存在词元与对象查询一同进入解码器，再由存在预测头得到$p_{\mathrm{pres}}$；查询$i$的局部分数$s_{i\mid\mathrm{pres}}$表达概念存在条件下的匹配，最终使用$s_i=p_{\mathrm{pres}}s_{i\mid\mathrm{pres}}$。训练分别提供存在标签与实例范围；局部分类损失在概念存在的图像上计算，使全局缺席主要由共享分支学习。不能仅用“图里有车”这一标签学出所有车辆的像素边界。

文字先取得词元特征，区域示例则同时编码框位置、示例中的视觉内容及正负标签，与文字共同形成概念条件。图像特征通过融合编码读取条件，对象查询再细化框与掩码；负示例告诉模型哪些相似对象不属于当前概念，不是将该区域的像素全部删除。训练记录包含概念、存在与否、全部匹配实例及提示。存在头用二元交叉熵覆盖正图与难负概念；概念存在时，实例按类别和范围代价唯一匹配，再用局部得分、框和掩码误差训练查询。缺席概念不向局部分类分支施加逐查询负梯度，职责与普通DETR的无对象训练有别。

这些接口来自有阶段区分的参数学习。SAM 3先用图文对预训练视觉和文字编码器，再以图像、短语和掩码三元组训练图像检测器及两侧编码器；三元组既有伪标注也有人工标注。高质量图像阶段加入存在头，并根据漏检采正区域、根据高分误检采负区域，累加示例以模拟纠错。该阶段联合降低框$L_1$、广义IoU、掩码焦点与Dice以及分类、存在误差。图像检查点因而已经学过概念与范围关联；使用时更换“校车”文字只改变条件，不重新学习权重。后续视频跟踪冻结主干的阶段不应倒推为前面的图像训练也一直冻结。

若教学中存在分数为0.9，两个实例局部分数为0.8和0.6，最终分数为0.72和0.54；不存在概念的存在分数为0.05时，同样局部分数只剩0.04和0.03。最终阈值决定保留哪些查询，再将各掩码放回原图并保持独立实例编号。共享分支的一个错误会同时改变所有查询，故不能只看保留掩码的边界而忽略被整体压低的概念。

对于不存在的概念，即使某个候选与“车”这一宽泛外观相似，共享存在判断也可以抑制错误输出。反过来，共享存在分支若把“校车”判为不存在，多个局部查询都可能被压低，形成相关的漏检。评价因此分别检查概念有无、每实例召回和掩码边界。简短概念与区域示例扩大了目标选择方式，但尚不足以推断任意关系嵌套或长指令都已解析；视频身份传播在\ref{chap:vision-video-recognition}章继续讨论。

\subsection{软边界与前景贡献的估计}
\label{sec:vision-recognition-matting}

发丝和半透明边缘的一个像素可能同时收到前景与背景的光。二值掩码只能选择保留或删除，无法表达这种混合。\term{图像抠图}（image matting）估计前景贡献比例$\alpha\in[0,1]$。在线性颜色空间的简化合成模型中，每个像素满足
\begin{equation}
 \bm{x}=\alpha\bm{f}+(1-\alpha)\bm{b}\text{，}
 \label{eq:vision-matting}
\end{equation}
其中$\bm{x},\bm{f},\bm{b}\in\mathbb{R}^3$是观测、前景与背景颜色。只有三个颜色观测，却有前景、背景和$\alpha$共七个未知标量，需要额外条件或先验；二值分割标签并不能唯一解出它们。

\subsubsection{三分图条件与alpha细化：Deep Image Matting}
\label{sec:vision-recognition-deep-matting}

Deep Image Matting使用\term{三分图}（trimap）指定确定前景、确定背景和未知区域，把RGB图与这一提示通道共同输入编码解码网络，再用细化网络改善$\alpha$边界。确定前景约束$\alpha=1$，确定背景约束$\alpha=0$，学习的重点是未知带中的混合比例\scite{vision-1703-03872}

训练可利用已知前景、背景与$\alpha$构造合成图像，监督$\alpha$误差，并按式\eqref{eq:vision-matting}检查预测贡献比例能否重建合成观测。使用时输出$\alpha$及其已知区约束；要把对象放到新背景$\bm{b}'$，仍需估计或取得前景颜色$\bm{f}$，再合成为$\alpha\bm{f}+(1-\alpha)\bm{b}'$。只取得$\alpha$并不意味着前景颜色已无污染。

RGB与三分图组成四通道输入，编码器压缩空间信息、解码器恢复未知带的$\hat\alpha$，细化网络读取图像及初步$\alpha$再预测细节。训练在未知区$U$中比较$|\hat\alpha_p-\alpha_p|$，并比较$\hat\alpha_p\bm f_p+(1-\hat\alpha_p)\bm b_p$与合成观测的颜色误差。两项通过同一混合公式连接：前者约束贡献比例，后者约束该比例的颜色后果。已有视觉权重可初始化编码层，新增输入和解码部分由合成监督学习，细化阶段再利用已训练的初步结果。三分图可以从目标$\alpha$构造未知带以训练，但使用时必须实际取得提示，不能把真实$\alpha$隐藏在输入中。

例如灰度前景为0.8、原背景为0.2、$\alpha=0.25$，观测为0.35。换成0.6的背景，正确合成为$0.25\times0.8+0.75\times0.6=0.65$；若二值化为背景，结果为0.6，发丝的贡献消失。若直接把观测0.35当成前景颜色，又会混入旧背景。提示误把发丝标成确定背景时，网络还可能被错误条件限制。边界、贡献比例和最终合成质量应分别核验。

\section{数量、姿态与空间测量}
\label{sec:vision-recognition-measurement}

“有几辆车”“人的手腕在哪里”“路牌有多远”分别要求数量、具有部位语义的位置以及带坐标或尺度的空间量。它们可以调用前面的对象范围，但监督和误差不能由检测指标替代。

要把手腕或路牌放到三维空间中，先看一个已经选定的表面点如何成像。\emph{针孔相机}（pinhole camera）用相机中心和一个成像平面描述投影：空间点与相机中心的连线穿过图像平面，交点就是它的像素位置。这里忽略镜头畸变，采用位于相机前方的虚拟成像平面，使像素坐标向右、向下增加，相机光轴向前。实际镜头存在畸变时，应先用标定参数校正，或在投影中加入畸变模型。

同一点可以在不同坐标系中表达。\emph{世界坐标}（world coordinates）是为整个场景选定的公共参照，\emph{相机坐标}（camera coordinates）则以当前相机中心为原点。记世界点为$\bm{P}_{\mathrm w}\in\mathbb{R}^3$，相机点为$\bm{P}_{\mathrm c}=(x_{\mathrm c},y_{\mathrm c},z)^\top$。\emph{外参}（extrinsic parameters）$\bm{R}\in\mathbb{R}^{3\times3}$和$\bm{t}\in\mathbb{R}^3$把前者变成后者，即$\bm{P}_{\mathrm c}=\bm{R}\bm{P}_{\mathrm w}+\bm{t}$，其中$\bm{R}^{\top}\bm{R}=\bm{I}$、$\det\bm{R}=1$。这里$\bm{t}$是世界到相机变换的平移，不是相机中心的世界坐标；相机中心为$\bm{C}_{\mathrm w}=-\bm{R}^{\top}\bm{t}$，逆变换为$\bm{P}_{\mathrm w}=\bm{R}^{\top}(\bm{P}_{\mathrm c}-\bm{t})$。

只知道相机点还不能得到像素，因为焦距和图像原点仍需确定。\emph{内参}（intrinsic parameters）描述这一步映射。忽略像素轴倾斜时，内参矩阵为
\[
 \bm{K}=\begin{pmatrix}
 f_x&0&c_x\\0&f_y&c_y\\0&0&1
 \end{pmatrix}\text{，}
\]
其中$f_x,f_y>0$是以像素为单位的焦距，$(c_x,c_y)$是主点，即光轴落到图像平面的位置。沿光轴深度$z>0$的点与投影平面构成相似三角形，横纵方向分别给出$u-c_x=f_xx_{\mathrm c}/z$、$v-c_y=f_yy_{\mathrm c}/z$。把这两个关系写在同一个齐次式中，得到
\begin{equation}
 z\begin{pmatrix}u\\v\\1\end{pmatrix}
 =\bm{K}(\bm{R}\bm{P}_{\mathrm w}+\bm{t})\text{。}
 \label{eq:vision-camera-projection}
\end{equation}
齐次坐标的末项让平移可以与投影共同表达，并不增加一个物理空间维度。$z=0$时除法未定义，$z<0$的点在相机后方，不能按当前可见图像点使用。

式\eqref{eq:vision-camera-projection}也说明像素为何不能单独确定深度。令$\tilde{\bm{p}}=(u,v,1)^\top$，则沿这条像素射线的相机点都是$\bm{P}_{\mathrm c}(z)=z\bm{K}^{-1}\tilde{\bm{p}}$，其中$z>0$待定。光轴深度与到相机中心的直线距离$\rho$不是同一个量：
\begin{equation}
 \rho=\lVert\bm{P}_{\mathrm c}\rVert_2
 =z\sqrt{1+\left(\frac{u-c_x}{f_x}\right)^2
            +\left(\frac{v-c_y}{f_y}\right)^2}\text{。}
 \label{eq:vision-depth-range-measurement}
\end{equation}
教学中取$f_x=f_y=800$像素、主点$(320,240)$和像素$(520,240)$。若$z=4$米，反投影得到$(1,0,4)$米，直线距离为$\sqrt{17}\approx4.123$米；若$z=8$米，点$(2,0,8)$仍投到同一像素。两次投影相同，测量量却不同。

多视图增加了射线约束，但还需区分\emph{相对尺度}（relative scale）与\emph{度量尺度}（metric scale）。前者只确定长度比例，后者还确定一单位对应多少米。若把全部世界点和相机平移同时乘$a>0$，相机点与$z$也乘$a$，式\eqref{eq:vision-camera-projection}的像素不变。已知相机基线、已知物体长度或带单位的深度观测可以确定这个比例；只把首相机设为原点，只固定了参照位置和朝向，没有测出尺度。网络从人体或街景训练分布预测的尺度是附加先验，也须在当前设备与场景中验证。

同一个世界点进入第二幅图时，要先回到世界坐标，再进入第二相机。对两相机外参$(\bm{R}_1,\bm{t}_1)$与$(\bm{R}_2,\bm{t}_2)$，消去世界点可得
\begin{equation}
 \begin{aligned}
 \bm{P}_{\mathrm c,2}
 &=\bm{R}_{2,1}\bm{P}_{\mathrm c,1}+\bm{t}_{2,1},\\
 \bm{R}_{2,1}&=\bm{R}_2\bm{R}_1^\top,\qquad
 \bm{t}_{2,1}=\bm{t}_2-\bm{R}_2\bm{R}_1^\top\bm{t}_1\text{。}
 \end{aligned}
 \label{eq:vision-relative-camera-measurement}
\end{equation}
将第一幅图的反投影点代入，再用第二幅图的$\bm{K}_2$投影，就得到一个可检查的跨图位置。裁剪和缩放也必须进入内参：若横纵缩放为$a_x,a_y$，裁剪原点为$(u_0,v_0)$，则新焦距为$(a_xf_x,a_yf_y)$，新主点为$(a_x(c_x-u_0),a_y(c_y-v_0))$。把新像素配旧内参会直接改变射线方向。

落在图像内且位于相机前方，只是可见的必要条件。\emph{遮挡}（occlusion）指同一射线上更近的表面挡住更远的表面；不透明场景中的像素通常记录最先遇到的表面。两幅图之间比较一个投影点时，还要检查源图深度是否与该点相容，否则取到的可能是挡在前面的车辆。人体被挡住的手腕可以由结构先验推断，但不是已经观察到的关节点。以下姿态、深度和场景恢复都沿用这些坐标、尺度与可见性约定；计数只需规定个体与区域规则，不必为此先重建整个场景。

\subsection{对象数量的估计}
\label{sec:vision-recognition-counting}

稀疏对象可以先逐个识别再计数，密集遮挡场景则可以直接估计数量的空间分布。前者保留对象清单，后者强调总量，两者取得正确总数的原因可能不同。

\subsubsection{检测实例汇总与重复计数控制}
\label{sec:vision-recognition-instance-counting}

给定去重后的检测集合，可按类别、区域与可见条件筛选，再把满足条件的指示量相加。若区域为$R$，以框中心$\bm{c}_i$是否在$R$内为计入规则，则
\begin{equation}
 \hat n_R=\sum_i
 \mathbf{1}\{c_i\in C_{\mathrm{target}}\}
 \mathbf{1}\{s_i\geq\eta\}
 \mathbf{1}\{\bm{c}_i\in R\}\text{。}
 \label{eq:vision-detection-count}
\end{equation}
这里$c_i$表示类别，$\bm{c}_i$表示二维中心，两者字形不同。中心落入与“任意可见部分落入”是两种不同规则，边界目标会得到不同计数。

例如真实有5辆车，漏掉1辆，另1辆被重复预测，总数仍可能为5。只看计数误差会掩盖两种错误；应同时检查对象匹配，尤其是遮挡和边界处的实例。跨帧统计还要避免同一辆车被多次计入，需由跟踪维持身份。

\subsubsection{多尺度密度预测与积分：MCNN}
\label{sec:vision-recognition-mcnn}

在人群过密、单个身体范围难以标注时，可用头部点标注学习\term{密度图}（density map），即每个位置分摊多少计数贡献。MCNN以多个不同感受范围的卷积分支处理尺度变化，连接各列特征后回归一张密度图\scite{vision-mcnn2016}

若标注点为$\bm{p}_i$，可用归一化核构造监督
\begin{equation}
 D(\bm{p})=\sum_{i=1}^{n}G_{\sigma_i}(\bm{p}-\bm{p}_i),
 \qquad
 \sum_{\bm{p}}G_{\sigma_i}(\bm{p}-\bm{p}_i)=1\text{。}
 \label{eq:vision-density-target}
\end{equation}
原工作令核宽与邻近头部点的平均距离相关，用局部拥挤程度近似尺度变化。卷积列和汇总层以预测密度与目标密度的平方误差训练。使用时对预测网格求和得到人数；密度可以是非整数分布，最终总数也可保留实值用于评价。

更具体地，若第$i$点最近$k$个邻点的平均距离为$\bar d_i$，可设$\sigma_i=\beta\bar d_i$，其中$\beta>0$及邻点数由标签构造约定。它只是从点间距得到尺度近似，孤立点或点数不足时需回退核宽。每列卷积读取不同邻域，共同输出单位为“该密度格分摊人数”的$\hat D_i$。经验目标可写为$L=(2n)^{-1}\sum_{i=1}^n\sum_p(\hat D_i(p)-D_i(p))^2$，每格梯度与$\hat D_i-D_i$成正比，训练将它传回汇总层及各列。原MCNN分别预训练卷积列后再联合适配；使用已训练参数时没有逐对象分类步骤。

若连续密度单位为人／像素面积，缩放$(x',y')=a(x,y)$使面积扩大$a^2$，需取$D'(x',y')=D(x'/a,y'/a)/a^2$，变量代换才给出$\int D'=\int D$。离散网格应相应重分配计数质量；输出步幅变大时，标签下采样也应保存总和。裁剪训练时先按规定中心规则确定哪些点计入，再为这些点在裁剪域构造归一化核；把原密度图任意裁一块，会把边外对象的核尾也计入，并丢掉边内对象的一部分质量。

三个点各贡献1，整图密度之和便为3。若一个边缘点的核有一半落到图外，直接截断后该点只贡献约0.5；有限网格构造标签时需按计入规则归一化有效核。图像缩放也要保持总和：将每个密度单元分成四个，四个新单元应分享原来的质量，不能每个都复制原值。

总和正确仍可能把本属左半边的人数移到右半边。应同时看局部计数、空间误差与场景变化。透视改变对象尺度，新的拍摄高度和密度范围可能超出训练条件；密度图没有自动恢复每个人的身份或完整位置。

\subsection{关节与姿态的估计}
\label{sec:vision-recognition-pose}

人体关节不同于图像中任意可匹配的角点：左腕、右肘具有固定部位含义。多人场景还必须知道这些部位属于谁。二维位置位于图像平面，三维姿态则要增加深度与坐标约定。

\subsubsection{二维部位定位与关联：OpenPose}
\label{sec:vision-recognition-openpose}

自顶向下路线先检测人，再在各人区域内定位关节；多人重叠时，候选框和裁剪质量影响后续姿态。OpenPose采用自下而上路线，在整图预测关节热图以及连接部位的\term{部位亲和场}（part affinity field, PAF），再组装人体\scite{vision-1611-08050}

热图在对应关节附近取得高响应，峰值形成候选位置。PAF是每种肢段的一张二维向量场：真实肢段附近写入从一种部位指向另一种部位的单位向量，其他位置为零，重叠肢段按规则合并。训练从关键点标注构造两类目标，对有效位置计算热图与向量场误差，并逐阶段细化。

设候选肘部、腕部为$\bm{a},\bm{b}$，方向$\bm{u}=(\bm{b}-\bm{a})/\lVert\bm{b}-\bm{a}\rVert_2$，对应PAF为$\bm{L}$。沿候选肢段平均方向一致程度：
\begin{equation}
 s(\bm{a},\bm{b})=
 \int_0^1
 \bm{L}((1-\lambda)\bm{a}+\lambda\bm{b})^\top
 \bm{u}\,\mathrm{d}\lambda\text{。}
 \label{eq:vision-paf-score}
\end{equation}
实现中在有限个点上采样求平均，再按部位连接类型作匹配并组合人体。$\bm{a}=\bm{b}$时方向未定义，应排除这一退化候选。

每类关节的目标热图可由标注位置周围的高斯响应构造，多人同类关节按规定合并；每种肢段在其附近带状区域写入方向，重叠区域平均有效方向。各阶段对有标注位置计算热图及PAF平方误差，未标注身体区域以有效掩码排除，梯度共同更新特征及分支。推理先从热图峰值取点，再对每种肢段建立候选两端的评分矩阵，禁止同一部位候选在该连接类型中任意占多条连接，按匹配结果逐步拼成身体。人体组合还需利用共享关节编号合并连接，不能把所有高分线段简单相加。

假设五个采样点的方向内积为$(1,1,0.8,1,0.7)$，平均为0.9；连到另一人的手腕时变为$(0.1,0,-0.2,0,0.1)$，平均为0。两个手腕都定位正确，也仍需这一关联步骤才可能把手臂接对。投影重合、严重遮挡和交叉肢段会使证据混杂，因此关节位置与人体组装应分开评价。

\subsubsection{二维关节到三维姿态：Martinez基线}
\label{sec:vision-recognition-lifting}

取得二维关节后，可以把按固定部位顺序排列的$2j$维坐标输入回归网络，输出$3j$维三维坐标。Martinez等人的基线使用全连接层与残差模块学习这一映射，使二维检测和三维恢复的误差来源能够分开研究\scite{vision-1705-03098}

训练规范化输入坐标，以对应的三维姿态作回归监督。三维目标通常相对于根关节表示，并遵守指定相机坐标、尺度和关节集合。使用时必须复现训练的规范化，再将输出还原为该协议中的坐标；关节顺序错位会使一个部位的输入被当成另一个部位，即使维度仍然匹配。

记按训练集统计规范化的二维坐标为$\bm{x}_i\in\mathbb{R}^{2j}$，相机坐标中的三维关节减去根关节后为$\bm{y}_i\in\mathbb{R}^{3j}$。提升网络先把$\bm{x}_i$投到较宽隐层，残差块以$\bm{h}'=\bm{h}+F_{\bm{\theta}}(\bm{h})$更新，其中$F$包含全连接、归一化、非线性及训练时随机屏蔽；最后线性层恢复三维输出。以$(1/n)\sum_i\lVert g_{\bm{\theta}}(\bm{x}_i)-\bm{y}_i\rVert_2^2$训练，从初始化权重逐批更新，保留网络和训练集规范化统计。输出逆规范化后得到相对根关节坐标，不含被减去的根位置；要得到人体到相机的距离，还需要另估根平移，不能直接由这组相对坐标读取。

节首教学像素$(520,240)$可对应$(1,0,4)$或$(2,0,8)$。人体长度先验能够在这些解释中形成偏好，但真正尺度来自哪个标注协议仍需说明。真实二维输入隔离关节检测误差，检测二维输入则要先逆人体裁剪并按同一规范化处理。已知相机的多视图关节可以通过射线求交或最小重投影误差估计三维点，增加的是观测，不是把单目回归网络的输出自动变成多视图证据。

向前或向后弯曲的肢段也可能有相近二维投影，网络用训练中的人体结构先验选择一种三维解释。评价要说明是否允许旋转、平移及尺度对齐：相似变换对齐后的关节误差较小，不意味着原始输出已具有正确朝向或距离。

\subsubsection{从关节坐标到可提示的人体表面恢复}
\label{sec:vision-recognition-sam3d-body}

两个三维骨架可以具有相同关节位置，却对应不同肩宽、躯干厚度和手指表面。若需要把人体轮廓投回图像，关节点还不够。\emph{人体网格恢复}（human mesh recovery, HMR）估计由顶点和连接面片组成的人体表面，同时解释其姿态、体形与相机。表面顶点数只是表示容量，不能证明被遮挡的背部已经恢复真实形状。

参数化人体表示用少量身体参数生成骨架和网格，使输出遵守同一人体连接结构。SAM 3D Body使用Momentum Human Rig（MHR），将骨架结构、关节姿态和表面形状分开控制。骨架结构规定连接和骨骼长度，姿态规定各关节如何相对旋转，表面形状规定骨架周围的身体外形。记这些参数为$\bm{\beta}_{\mathrm{sk}},\bm{\theta}_{\mathrm{pose}},\bm{\beta}_{\mathrm{surf}}$，已提供的人体模型分别产生关节和顶点：
\[
 \bm{J}=\mathcal{J}(\bm{\beta}_{\mathrm{sk}},
                   \bm{\theta}_{\mathrm{pose}}),\qquad
 \bm{V}=\mathcal{V}(\bm{\beta}_{\mathrm{sk}},
                   \bm{\theta}_{\mathrm{pose}},
                   \bm{\beta}_{\mathrm{surf}})\text{。}
\]
这里$\bm{J}\in\mathbb{R}^{j\times3}$、$\bm{V}\in\mathbb{R}^{v\times3}$，$j,v$由模型资产规定；这两个表达式只说明输出分工，不是把MHR的全部变形计算替换成一个线性式。把关节或顶点经相机投影，才得到可与原图比对的位置和轮廓\scite{vision-sam3d-body-2026}。

模型首先取得人体裁剪及裁剪到原图的变换，归一化后由共享图像编码器产生空间特征。已发布版本可使用DINOv3视觉主干，身体和手部解码器由人体任务监督学习；推理还需加载相应MHR模型资产，只有网络权重不足以确定顶点连接和变形规则。身体解码器的查询中包含初始MHR和相机估计、辅助二维和三维关节点，以及可选的用户提示。二维关键点提示写成$(x,y,\ell)$，其中$\ell$指出关节身份；位置编码与学习的关节嵌入使“左腕在这里”不同于一个无部位含义的点。掩码提示经卷积编码后加到图像特征，提供人体占据区域。查询读取图像和提示，输出姿态、表面形状、骨架及相机参数，再通过MHR资产恢复完整网格。

手部占据整人裁剪中的像素可能很少，模型另设可选手部裁剪和独立手解码器，仍共用图像编码器权重。身体解码器本已能输出包含手足的完整人体，手解码器提供手部专门训练形成的补充。原方法不把局部手参数直接替换进去后就结束：直接替换可能连带改变肘部，因而用手解码得到的腕位置和身体解码得到的肘位置再次提示身体分支，细化后按人体连接树合并。遮挡手部若没有可靠检测，应保留身体输出或关闭局部细化，而不是把不存在的裁剪当成有效观察。

训练样本的三维依据来自多阶段标注：人工关键点、可微拟合、多视图几何和密集关键点检测以不同组合产生监督，还包含合成和多视图资料。原论文首稿为2026年2月，并不据此认定系统首次发布也发生在该月。训练将二维关节放在相应裁剪坐标中监督，将身体三维关节减去骨盆、手关节减去手腕后比较；已标注手部才参与对应目标。关节使用带学习不确定性权重的$L_1$项，姿态与形状参数使用$L_2$项，另加关节限位和手框误差。随机多轮提示模拟纠错，对用户给定关节提高监督权重，部分三维项逐步升权。训练得到的是图像到人体参数的映射，不是为每张新图重新收集真实三维网格。

教学中设一个人左腕被遮住，无提示预测投回原图$(120,150)$，而用户根据额外依据给点$(128,146)$。两者像素残差为$\sqrt{8^2+4^2}\approx8.94$像素。加入关节提示后应重跑解码、恢复网格并重新投影，检查该残差以及肘腕连接、手部形状和身体轮廓是否共同改变；加入人体区域掩码时则主要检查轮廓条件。这个计算只是定义需要检查的量，没有假造提示后的模型误差或宣称误差必然下降。

使用时人体框可由外部检测器取得，框误差属于完整流程的一部分。用户点需先变到人体裁剪坐标，预测关节和网格在相机条件下投影后，再用裁剪逆变换回原图。提示轮次可由用户确认或预算停止，最终保存人体编号、MHR参数、相机、网格及提示来源。检查点的量纲与首相机约定需要保留；未知真实身高时，不能把模板的身体尺寸当成该人的测量结果。

一个轮廓相符的网格仍可能把手腕置于不同深度，节首的尺度歧义也仍存在。应分别检查关节、顶点表面、手足细节、提示一致性和相机，而不只比较一张渲染图。相机坐标中的合理人体不能直接给出世界坐标的真实身高、动作意图或生物力学量；后两类还需时序、力学与采集条件。连续姿态及身份关联由\ref{chap:vision-video-recognition}章承接。

\subsection{深度与距离的估计}
\label{sec:vision-recognition-depth}

现在把测量对象从关节换成每个可见表面像素的光轴深度。节首的式\eqref{eq:vision-camera-projection}已经给出像素射线，尚未确定的是射线上的$z$。单幅外观可以提供训练先验，多幅已知相机的图像可以提供跨射线约束；两条路线都需说明输出采用深度、逆深度还是其他相对量，以及是否已经确定度量尺度。

\subsubsection{单幅相对深度与度量适配：Depth Anything V2}
\label{sec:vision-recognition-depth-anything}

单目模型利用物体形状、透视、遮挡和训练先验，从一张图预测密集深度。Depth Anything V2以合成数据的精确深度训练较强教师，再为大量真实图像产生深度伪标签，供不同规模学生学习；图像编码器与稠密解码器把表示接回各空间位置\scite{vision-2406-09414}

这一训练桥接应区分两个信息来源：合成图像提供几何标签，真实图像提供目标外观，伪标签质量依赖教师。原工作以已有DINOv2编码器和密集预测变换器（dense prediction transformer, DPT）解码器作为起点；后者把不同层空间特征逐步融合、上采样成逐位置数值，不把全图汇总向量当作一张深度图。较大教师先用合成RGB和精确深度训练，教师为真实RGB产生伪深度后，学生只从这些真实图与伪标签对继续学习，教师标签不随学生每次梯度更新而变成真实测量。

基础检查点预测的是尺度、平移不变的相对逆深度。为看清“允许尺度和平移”如何作用，设预测相对量为$\hat q_p$、目标为$q_p$，一个教学用的平方对齐目标可以写成
\begin{equation}
 \min_{a,b}\frac1{|\Omega|}
       \sum_{p\in\Omega}(a\hat q_p+b-q_p)^2\text{。}
 \label{eq:vision-relative-depth-alignment}
\end{equation}
这里$\Omega$是当前图的有效监督位置，$a,b$是比较时的对齐量，不是已知米制相机参数。对$a,b$求导为零得到$\sum\hat q_p(a\hat q_p+b-q_p)=0$、$\sum(a\hat q_p+b-q_p)=0$；若预测有非零方差，便有$a=\operatorname{Cov}(\hat q,q)/\operatorname{Var}(\hat q)$、$b=\bar q-a\bar{\hat q}$。若全图预测为常数，尺度无从估计。这个平方例子解释对齐消去什么信息，不冒称原论文尺度平移不变目标的逐项实现。原训练还比较相邻位置残差的梯度，防止整体对齐虽合适、局部边界却过于平滑。

例如预测$(1,2,3)$与目标$(3,5,7)$在$a=2,b=1$下完全重合；对齐误差为零没有证明原预测的数值已是正确逆米单位。对伪标签样本，原工作忽略损失最大的部分区域以降低噪声影响，并加入特征对齐以保持已有语义信息；这不保证被忽略位置都是真噪声。训练的编码器和解码器分别按设定学习率逐批更新，验证选择可用检查点，不在使用时为每幅图再求真实深度标签。伪标签学习的一般依据可参考\BookChapterRef{chap:data-derived-knowledge-acquisition}{卷八“从数据中获得知识”章}，本节的新增条件是深度约定与空间有效区域。

使用相对检查点时，将原图按模型尺寸和归一化处理，编码成空间特征，解码出$\hat q$，再把浮点图插值回原图有效区域并保留相对逆深度身份。颜色图只是显示映射，不能从色值直接读米。度量适配改用带单位的深度标签更新已有网络，使输出在该数据、最大深度及相机条件下对应实际距离；模型版本与量纲需随结果一同保留。

若三个位置的相对逆深度为$(0.8,0.4,0.2)$，可以按该约定判断第一个较近，却不能未经校准说深度为$(1.25,2.5,5)$米。用带实际单位的深度标签适配后，模型才学习特定数据条件下的度量映射；即便如此，相机和场景变化仍可能带来尺度误差。边界位置、前后排序和实际距离应分别评价。

\subsubsection{相机条件下的代价体与深度推断：MVSNet}
\label{sec:vision-recognition-mvsnet}

如果有多张重叠照片和已知相机，同一表面在不同图像中的对应可以提供深度约束。MVSNet选定参考图，对每个候选深度$d_k$，把参考像素反投影成空间点，再投到其他相机，在源图特征中插值取值。整个平面上的映射可写为依赖$d_k$的可微单应变换\scite{vision-1804-02505}

这一步可由节首投影直接算出。参考像素$\tilde{\bm p}_r$在假设深度$d_k$下为$d_k\bm K_r^{-1}\tilde{\bm p}_r$，进入源相机得到$\bm R_{s,r}d_k\bm K_r^{-1}\tilde{\bm p}_r+\bm t_{s,r}$。由于$\bm e_3^\top\bm K_r^{-1}\tilde{\bm p}_r=1$，其中$\bm e_3=(0,0,1)^\top$，除去齐次比例$d_k$便得到
\begin{equation}
 \begin{aligned}
 \tilde{\bm p}_s&\sim\bm H_{s,r}(d_k)\tilde{\bm p}_r,\\
 \bm H_{s,r}(d_k)=
 \bm K_s\left(\bm R_{s,r}
 +\frac{\bm t_{s,r}\bm e_3^\top}{d_k}\right)\bm K_r^{-1}\text{。}
 \end{aligned}
 \label{eq:vision-mvsnet-homography-process}
\end{equation}
该式只针对参考相机中$z=d_k$的平面，不能用一个深度平面代替整幅起伏场景。特征图尺度不同于图像时，先按步幅调整内参和采样坐标，再用双线性插值读取源特征；越界或后方投影不得当成有效观测。

正确深度应使同一表面的特征比较一致。将参考特征和变换后的各源特征记为$\bm{f}_r$，共有$n_{\mathrm{v}}$个视图，可以用方差形成匹配代价：
\begin{equation}
 \begin{aligned}
 \bar{\bm{f}}&=\frac1{n_{\mathrm{v}}}\sum_{r=1}^{n_{\mathrm{v}}}\bm{f}_r,\\
 \bm{c}(d_k,u,v)&=\frac1{n_{\mathrm{v}}}\sum_{r=1}^{n_{\mathrm{v}}}
   (\bm{f}_r-\bar{\bm{f}})\odot(\bm{f}_r-\bar{\bm{f}})\text{。}
 \end{aligned}
 \label{eq:vision-mvs-cost}
\end{equation}
这里$\odot$表示逐元素相乘。所有像素和候选深度的代价共同形成\term{代价体}（cost volume），它保留“每个位置在各个深度假设下多不一致”的信息。三维卷积在空间与深度轴上整理代价，经softmax得到深度权重$\pi_k$，用$\hat d=\sum_k\pi_kd_k$回归深度，并可作后续细化。训练只对具有有效参考深度的位置计算误差。

同一像素的教学标量特征在正确假设下为$(1,1.1,0.9)$，均值1、方差约0.0067；错误假设读到$(1,2,3)$，方差为$2/3$。真实网络在多通道、多位置及深度轴上整理代价，不能仅靠这一个方差保证深度正确。若候选深度为2、4、6米，权重为0.1、0.7、0.2，期望回归得到4.2米，不必等于某个离散候选。双峰权重也会把结果放在两个表面之间，因而概率集中度和有效区域需要另行检查。

训练记录含参考图、若干重叠源图、已知相机及有效参考深度。先确定深度范围和候选间隔，形成特征代价体和预测深度，在有效像素上以$L_1$比较初步及细化结果，再更新二维特征提取、三维代价整理及细化网络。使用时保留这些训练参数，但相机是当前图集提供的条件，不是从训练集复制的外参。恢复深度到参考网格后，可按置信和多视图几何一致性过滤，再反投影合并点云；每幅图看到的表面不同，直接叠加所有点仍可能留下遮挡冲突。若候选数为$k$、特征网格为$h'\times w'$、通道为$d$，仅代价体就需$O(kh'w'd)$存储。

在平行、已校正的双目简化情形中，焦距$f$、相机基线$b$和水平视差$\Delta u$满足$d=fb/\Delta u$。取$f=800$像素、$b=0.2$米、$\Delta u=40$像素，得到$d=4$米。视差误差1像素在此处约引起0.1米深度误差；同样误差在更远处会被放大。这一关系说明已知基线如何提供实际尺度。

图\ref{fig:vision-depth-hypotheses}展示两个深度假设如何查询不同源图位置。弱纹理可能让多个候选都具有相近代价，遮挡使源图看到另一表面，反光破坏外观一致性，相机误差又使正确深度也无法对齐。需要同时报告深度误差、可恢复面积与候选数带来的存储成本。

\begin{figure}[htbp]
 \centering
 % 针孔投影的几何关系；示意深度候选，不按物理长度绘制。
\begin{tikzpicture}[font=\figurefont\small,x=1cm,y=1cm,
 ray/.style={draw=BookBlue,line width=0.9pt},
 alt/.style={draw=BookAmber,dashed,line width=0.9pt}]
 \coordinate (a) at (0,0);
 \coordinate (b) at (7,0);
 \coordinate (p) at (2,3);
 \coordinate (q) at (3,4.5);
 \draw[ray] (a)--(q);
 \draw[ray] (b)--(p);
 \draw[alt] (b)--(q);
 \draw[draw=BookInk,line width=1.2pt] (-0.3,1)--(2.2,1);
 \draw[draw=BookInk,line width=1.2pt] (4.5,1)--(7.5,1);
 \fill[BookInk] (a) circle(1.1mm);
 \fill[BookInk] (b) circle(1.1mm);
 \fill[BookTeal] (p) circle(1.1mm);
 \draw[BookAmber,fill=white,line width=0.9pt] (q) circle(1.1mm);
 \fill[BookBlue] (0.667,1) circle(1mm);
 \fill[BookBlue] (5.333,1) circle(1mm);
 \draw[BookAmber,fill=white,line width=0.9pt] (6.111,1) circle(1mm);
 \node[below] at (a) {参考相机};
 \node[below] at (b) {源相机};
 \node[above right] at (0.667,1) {同一参考像素};
 \node[below] at (5.333,0.9) {$u'_1$};
 \node[below] at (6.111,0.9) {$u'_2$};
 \node[left] at (p) {深度$d_1$};
 \node[right] at (q) {深度$d_2$};
 \node[right] at (7.5,1) {源图平面};
 \draw[draw=BookMuted,<->,line width=0.7pt] (0,-0.8)--(7,-0.8)
   node[midway,below] {已知相机相对位置};
\end{tikzpicture}

 \caption{已知相机下，参考像素沿射线形成不同深度假设，分别投到源图的不同位置。图为针孔几何教学示意，只表示对应关系，不按物理长度绘制。跨图比较可排除不一致的假设，但遮挡或重复纹理会使判断不唯一。}
 \label{fig:vision-depth-hypotheses}
\end{figure}

\subsection{相机、场景几何与空间语义的恢复}
\label{sec:vision-recognition-geometry}

前一方法把相机作为已知条件。若相机未知，需要共同解释视图关系、表面点和投影。经典流程先寻找对应，再估计相机和三维点；学习式方法可以直接预测这些几何量，仍需处理共同坐标与尺度。

\Needspace{65mm}
\subsubsection{成对点图与全局对齐：DUSt3R}
\label{sec:vision-recognition-dust3r}

DUSt3R输入图像对，以两侧编码及跨图解码交换信息，为两幅图的像素分别输出三维\term{点图}（point map）：每个有效像素对应一个三维向量。两张点图都用第一幅相机的坐标表达，并预测每点置信度，使无需预先给出相机的成对几何成为直接输出\scite{vision-2312-14132}

训练记录含同一场景的图像对，以及由深度、内参和外参得到的真实点图与有效位置。两幅图共享视觉编码器，两个解码分支反复交换跨图信息，再分别输出各图像素的点图；第二张点图也变换到第一相机。编码解码权重可以从CroCo跨视图预训练初始化，输出头由三维监督学习，不把“使用时不输入相机”误写成“训练不曾使用几何标签”。

原工作分别用有效三维点到原点的平均距离规范化预测和目标，得到尺度无关点误差$e_i$。置信目标为$\sum_i(c_i e_i-\alpha\log c_i)$，其中$\alpha>0$，原实现以$c_i=1+\exp h_i>1$产生置信度。对固定$e_i>0$，该项导数为$e_i-\alpha/c_i$，较大误差倾向于较小置信，但受$c_i>1$限制；对数项使提高置信也需承担误差代价。以$\alpha=1$、教学误差$e=0.1$为例，无约束驻点为$c=10$；$e=2$的驻点0.5不在允许范围，便倾向于下边界。这是可靠性权重的训练关系，不是校准正确率。所有有效点误差反向更新网络，无效天空或透射区域按标签条件排除。

三张建筑照片形成多对输入后，每一对以自己的首图为参考，不能直接把全部点拼成场景。全局对齐优化各对点图到共同坐标的旋转、平移和尺度，使共享视图及可信点彼此相容。相机、深度与对应再从已对齐的几何联系恢复；一次成对前馈与多图优化是两个不同计算阶段。

先将有足够重叠的图对组成连接图，边$e$保存两张在同一对坐标中的点图$\bm X^{v,e}$及置信。对每幅视图求共同点图$\bm\chi^v$，对每条边求尺度$s_e>0$与刚体变换$(\bm R_e,\bm t_e)$，使$\sum_{e,v,p}c_p^{v,e}\lVert\bm\chi_p^v-s_e(\bm R_e\bm X_p^{v,e}+\bm t_e)\rVert_2$较小。相同边上的两图共用一个变换；对重复视图的每个像素共享同一$\bm\chi$，才把不同边联系起来。固定公共参照及尺度规范防止所有点任意缩到零。由初步连接初始化这些场景变量，按残差迭代优化至预设步数或改变量足够小，网络权重不随该场景对齐重新训练。

相机恢复再检查点图是否服从针孔关系。例如假设主点在中心、像素轴正交且焦距相同，可由$(u-c_x,v-c_y)\approx f(x/z,y/z)$对高置信点拟合$f$；对应三维点与二维像素可进一步估计相机姿态。将共同点图改写为相机、深度的反投影，还可在对齐目标中同时优化这些量。点图头本身没有硬性保证每个点满足同一个相机，所以相机恢复与跨图核验不能省略；连接图不连通时，也没有依据决定两个子场景的相对位置。

若甲乙图预测建筑宽度为2个单位，乙丙图为5个单位，需要先估计约2.5的相对尺度并对齐朝向位置。单对看似合理不保证多对闭环一致；弱重叠、动态行人和错误高置信点都会干扰对齐。没有额外度量依据时，共同场景仍只具有约定尺度。

\subsubsection{多视图共享特征与联合几何输出：VGGT}
\label{sec:vision-recognition-vggt}

逐对预测再优化可能重复处理同一视图。VGGT让多幅图的词元共同参与计算，交替进行帧内注意力与跨全部视图的全局注意力。帧内计算保留各图的结构，全局计算交换视角证据，随后由不同预测头输出相机、深度和点图；点跟踪头还接收参考图中的二维查询点\scite{vision-2503-11651}

训练使用对应几何量监督各预测头，把相关任务接在共享表示上。输出以首相机为参考，并沿用训练的尺度规范；这使多视图结果可在同一约定下比较，却没有自动赋予米单位。相机头给出的内外参和深度头的输出可以反投影成点，再与直接点图比较，从而检查联合输出是否几何一致。

一条训练记录是同一场景的多张RGB、每图相机、有效深度和共同坐标点图，跟踪监督还包含查询点的跨图位置及可见性。这些几何来自传感器、合成引擎或运动恢复结构，而不是只有图像类别。训练先把全部量变到首相机坐标，再以有效三维点到原点的平均距离$s$除相机平移、点坐标和深度。$s$是该场景的规范化长度，不作为使用阶段已知的真实米数。与DUSt3R在目标比较时再归一化预测不同，VGGT要求网络直接学习这套规范输出。

编码器用已有DINOv2视觉权重初始化，将每图变成带位置的图像块词元，并附一个相机词元和若干辅助寄存词元；后者是网络中暂存汇总信息的计算位置，不是实际相机或空间点。首图使用与其他图不同的学习词元标记参照。帧内层让图像词元与自己的相机词元交换信息，全局层再跨全部视图读取，因此相机词元始终与所属图像关联。最终相机头读取相机词元，稠密头将多层图像特征恢复为空间网格，分别产生深度、点图、置信及跟踪特征。

相机头的九维编码包含四维旋转四元数、三维平移和横纵视场角。四元数是表示旋转的四个归一化数，解码成满足正交条件的$\bm R_i$；视场角$\varphi_x,\varphi_y$在针孔与中心主点条件下给出$f_x=w/[2\tan(\varphi_x/2)]$、$f_y=h/[2\tan(\varphi_y/2)]$。因此得到的是当前网络图像网格的内参，而非不依赖缩放的固定焦距。深度$D_i$仍是第$i$相机自身光轴深度；只有点图和相机外参共同用首相机公共参照。恢复公共点可计算
\begin{equation}
 \bm{P}_{1,i}(u,v)=
 \bm R_i^\top\left[
 D_i(u,v)\bm K_i^{-1}(u,v,1)^\top-\bm t_i\right]\text{。}
 \label{eq:vision-vggt-unprojection-process}
\end{equation}
这里$\bm R_i,\bm t_i$沿用世界到第$i$相机约定，公共世界就是首相机。先将相机、像素和深度放在同一处理网格，才能逆投影；若改回原图坐标，也须相应改回$\bm K_i$。

训练以Huber误差比较相机编码，以置信加权位置误差和空间梯度误差监督有效深度与点图，以点位置误差和二元可见性目标监督跨图查询。Huber项在小残差区使用平方、大残差区使用线性代价，减轻极大相机残差的单项贡献，但不是自动排除错误标签。原模型联合更新视觉主干、交替注意力和任务头，按场景采样2至24幅图，图像与几何网格共同缩放；主干权重来自预训练，几何输出仍需这一步联合监督。原文使用有限迭代预算与AdamW，没有给任意图集恢复正确几何的保证。

使用三张建筑图时，加载既有检查点，按相同图像处理保留缩放和图号，一次联合前向得到每图相机和稠密量。可用式\eqref{eq:vision-vggt-unprojection-process}从深度构造公共点，也可读取直接点图；两条输出不应未经核验混拼。教学中首相机为单位旋转、零平移，主点$(320,240)$、焦距800、像素$(520,240)$与深度4个规范单位恢复点$(1,0,4)$，仍不能将4个单位直接写成4米。按置信和可见条件保留点，投回其他相机检查像素和深度是否相容。

跟踪接口需先把原图查询点变到当前处理网格，在查询图特征中插值取其描述，再与各图跟踪特征相关，估计跨图二维位置和可见性。输出位置逆变换回各自原图，并保留查询编号、源图及不可见状态。它可用于无时序图集中的对应，不是未给查询时自动生成一份车辆身份轨迹。训练中随机点的监督与使用时指定的窗角查询对应同一种位置输出，新增动态运动条件则由视频章解释。

多图前馈输出之后仍可进行几何优化，但应把这部分额外时间单独记录。图像数和每图词元数共同增加全局注意力成本，动态对象、超出训练范围的相机以及缺少重叠的图像也会削弱预测依据。新视角渲染将在\ref{sec:vision-generation-novel-view}节使用这些几何条件，生成外观与恢复可测量表面仍是不同目标。

\subsubsection{大规模离线几何与测试时场景压缩}
\label{sec:vision-recognition-vgg-t3}

建筑的视图从三张增到数百张，完整图集已可用，但同时交换全部词元变得昂贵。若每图有$p$个词元、共有$n_{\mathrm v}$幅图，则全局注意力比较$(n_{\mathrm v}p)^2$个查询键组合，键和值本身占$O(n_{\mathrm v}pd)$空间。高效注意力实现可以避免完整保存得分矩阵，却没有消除这些组合的计算。离线条件允许共同利用全部视图，需要改变的是场景信息的读取方式。

VGG-T$^3$把每个全局层中长度随视图增加的键值集合压缩到一个固定大小的多层感知机（multi-layer perceptron, MLP）。这里MLP将一个键向量映射成对应的值向量，用其参数保存当前场景的映射；它不是一个固定长度的三维点清单。原方法以VGGT权重初始化，保留图像编码、帧内计算和预测头，适配全局层的投影及新增测试时训练模块；查询键改用$L_2$归一化。值中的图像块部分先在各图网格上经过短二维卷积混合邻域，得到$\bm{v}'_i$，相机和辅助寄存词元不经过这项空间卷积。邻域混合使拟合目标带有局部上下文，避免仅复现同一词元的两个线性投影\scite{vision-vgg-t3-2026}。

对一个全局层，设当前图集产生$N=n_{\mathrm v}p$个键$\bm{k}_i$、查询$\bm{q}_i$和目标值$\bm{v}'_i$，维数均为$d$。测试时训练（test-time training, TTT）指在使用当前输入时临时更新一小组场景权重$\bm{\phi}$；预训练主干并未因这个名称变成未训练网络。为使键读出相应值，内层目标及查询计算可写为
\begin{equation}
 \begin{aligned}
 \bm{\phi}_*&\approx\argmin_{\bm{\phi}}
   \frac1N\sum_{i=1}^{N}
   \lVert T_{\bm{\phi}}(\bm{k}_i)-\bm{v}'_i\rVert_2^2,\\
 \bm{o}_i&=T_{\bm{\phi}_*}(\bm{q}_i)\text{。}
 \end{aligned}
 \label{eq:vision-vgg-t3-fast-weights}
\end{equation}
近似号强调实际只做规定次数的更新，不求一个已证明存在的精确最优映射。每层先从训练得到的初始化取得快速权重，按键值残差计算内层梯度，再将查询送入更新后的MLP；所得词元进入下一层，全部层完成后才由相机、深度和点图头解码。原文内层采用Muon优化，训练内层用一步，通常测试用两步；这属于所报告配置，不是任意场景的收敛结论。

平方残差目标还允许分批拟合全集。固定一次更新开始的$\bm{\phi}$，把词元划成批次$B_s$，则
\[
 \nabla_{\bm{\phi}}L=
 \frac1N\sum_s\sum_{i\in B_s}
 \nabla_{\bm{\phi}}
 \lVert T_{\bm{\phi}}(\bm{k}_i)-\bm{v}'_i\rVert_2^2\text{。}
\]
这一等式来自有限和的微分线性性，要求各批都在同一组权重上算梯度，累加后再更新；逐批改权重不是同一次全集梯度。可把未处理批放在主机内存，或跨设备分摊，降低单卡同时保存图集特征的需求。拟合和读取在固定MLP容量、分辨率及更新步数下随词元数线性增长，原图、编码特征、输出点图和数据传输仍随视图增多。

用标量线性模型$T_\phi(k)=\phi k$作教学简化，两个键值样本为$(1,2)$、$(2,4)$。在$\phi=0$处，平均平方目标的梯度为$-10$，取教学步长0.1更新到$\phi=1$，再更新一次到1.5；查询1.5读出2.25，还没有达到目标映射给出的3。真实方法用非线性MLP和不同优化器，这个例子仅说明“有限次拟合”与“精确保存全部键值”有差别。新视图若作为定位查询，可冻结已有各层快速权重只做读取，但原文相机头还保留建图图像的相机词元，不能把全部定位存储都说成与图集无关。

比较同一建筑逐步增加的视图，应同时记录内层步数、点图误差、相机误差、重投影、宽基线失败、总显存及总时间。原文在其基准中还指出相机预测较不精确，并采用直接点图头而非相机深度反投影评估点图。固定容量映射产生近似误差，不能由资源增长方式推断几何精度保持不变。这里完整图集先到齐；下一小节改变的是信息到达条件。

\subsubsection{流式几何状态、地图记忆与长期漂移}
\label{sec:vision-recognition-infinitevggt}

沿走廊前进时，第$t$幅图已经到达，返回走廊的画面却尚未发生。当前输出不能借未来视图纠正自身。\emph{因果几何预测}（causal geometry prediction）限制当前帧只读取自身与已到达历史。直接保留每层的历史键值，可以复用旧特征，但缓存随$t$增长；只保留最近几帧，又可能删去初始参照与回访线索。

InfiniteVGGT在已有因果几何模型上采用免训练的滚动缓存策略。图像仍经帧内和因果跨帧计算，输出相机、深度、点图及跟踪特征；变化在于哪些历史键值继续可读，不是为当前流重新拟合一个VGG-T$^3$式MLP。各层、各注意力头把第一帧的键值固定为锚点，因为几何输出以首相机为公共参照；其余历史与新到词元构成可裁减候选集\scite{vision-infinitevggt-2026}。

相邻画面含大量重复特征，原方法用归一化键的方向差异选择保留内容，避免读取完整注意力矩阵。对一个层和头，令候选非零键的单位向量为$\hat{\bm{k}}_i$，平均向量为$\bm{\mu}=m^{-1}\sum_i\hat{\bm{k}}_i$，多样性评分为
\begin{equation}
 s_i=-\frac{\hat{\bm{k}}_i^\top\bm{\mu}}
                 {\lVert\bm{\mu}\rVert_2}\text{。}
 \label{eq:vision-infinitevggt-diversity}
\end{equation}
与平均方向差别较大的键取得较大分数，并与自己的值一起保留。这个评分近似识别特征冗余，不是直接测量三维点的真实信息量；平均向量为零时余弦无定义，复现需有数值稳定和退化处理，不能从理想公式推出每组输入都可直接计算。

不同层的候选多样性不同。原文将层平均评分$\bar s_\ell$经softmax变成预算比例$b_\ell=\exp(\bar s_\ell/\tau)/\sum_j\exp(\bar s_j/\tau)$，用总可变预算$B$分配各层保留数，再按分数选前若干键值，最后加回不可裁减的首帧锚点。实现需把预算变为整数并核对总量；预算$B$、温度$\tau>0$和锚点成本均须记录。追加、评分、裁减、保留是每次到帧后的状态更新，权重本身不因裁减而接受梯度训练。

教学中令可变候选键方向为$(1,0)$、$(1,0)$、$(0,1)$，平均为$(2/3,1/3)$。前两者与平均方向余弦为$2/\sqrt5\approx0.894$，第三者为$1/\sqrt5\approx0.447$，第三者的负余弦评分更高；预算只容两项时，它和一个重复方向可以留下。示例说明了为何裁减可以保留不同方向，却没有证明被删方向在之后回访时不重要。

固定缓存控制的是工作状态，不是全部地图输出的永久存储。每帧仍产生点图，若全部写盘，总输出量仍随时间增加；当前帧到缓存的读取及评分也占计算。走廊返回起点时，应比较短程相机误差、累计轨迹漂移、回访位置重合与缓存成本。保留首帧和部分旧键有助于维持参照，却不是显式回环检测与全图校正。Long3D约万帧的有限序列可提供长程测试，不能证明无限时间稳定；地图参照与视频中的对象身份也不是同一个状态。

\subsubsection{三维几何与开放词汇语义的共同恢复}
\label{sec:vision-recognition-ff3r}

房间两幅照片都标出“椅子”，还不能知道它们是否指向同一张椅子的同一表面。二维语义可借投影进入三维，但几何错位会把椅面特征放到墙上；直接合并近邻三维点又可能把相邻椅子和背景混在一起。共同恢复需要分别保留位置、外观和可查询语义，并检查三者指向相同表面。

FF3R输入至少两幅多视图图像，沿用VGGT式几何词元，同时以冻结的CLIP-LSeg编码器取得逐位置语义特征。末层几何词元作为查询读取语义词元的键和值，使几何解码获得语义上下文。解码器输出深度和相机，并预测像素对齐的三维高斯：其中心由深度反投影确定，协方差描述空间扩展，颜色参数描述外观，不透明度描述遮挡贡献，另附语义向量$\bm{f}_i$。高斯是可投影、可聚合的场景表示，不应把其中心直接当作精确物体表面\scite{vision-ff3r-2026}。

为用图像监督这些三维属性，需要把它们投回观察相机。每个高斯中心用节首投影得到二维位置，空间协方差经投影的局部线性近似成为像平面的扩展区域；按深度先后计算每项对像素的贡献。若第$i$项在该像素的覆盖与不透明度合成系数为$a_i\in[0,1]$，前方透过比例为$T_i=\prod_{j<i}(1-a_j)$，则颜色和特征分别按$\sum_iT_i a_i\bm c_i$与$\sum_iT_i a_i\bm f_i$汇总，并按规定处理未覆盖背景。两者共用贡献权重，因此特征的监督也会反向影响所预测几何；完整新视角渲染在生成章展开，此处只用投回观测图的残差解释训练依据。

训练先从同一场景采样上下文视图，保留图号并采用相同裁剪尺度，初始化几何主干和深度头于已有VGGT权重，新增融合与输出部分随机初始化；几何交替注意力参与适配，CLIP-LSeg保持冻结。将预测三维表示投回观测视图，比较渲染RGB与真实RGB，以及渲染特征与CLIP-LSeg特征。原论文还明确使用VGGT给出的相机、深度伪标签作为蒸馏约束，深度只在高置信位置比较。因此“无需人工相机、深度和语义标注”不等于“只有RGB残差”，也不等于没有预训练监督。

若两视图看到同一表面，它们的语义特征还应在几何对应下相容。使用式\eqref{eq:vision-relative-camera-measurement}，将目标图像像素$\tilde{\bm{p}}_t$在预测深度$D_t$下反投影，再投到上下文图$c$：
\begin{equation}
 \bm{p}_c=
 \Pi\!\left[
 \bm{K}_c\left(
 \bm{R}_{c,t}D_t(\bm{p}_t)\bm{K}_t^{-1}\tilde{\bm{p}}_t
 +\bm{t}_{c,t}\right)\right]\text{，}
 \label{eq:vision-ff3r-warp}
\end{equation}
其中$\Pi$用第三坐标除前两坐标。按$\bm{p}_c$插值采样源特征，计算$1-\cos(\bm{f}_t,\bm{f}_{c\to t})$，只在投影落入图像且深度相容的位置求平均，再交换两图求另一方向。边界和遮挡筛选让“不曾看见同一表面”的位置不被强行拉到一起；特征不为零是余弦计算条件。

视图密集时，可按三维位置把高斯分到有限体素，即固定大小的立方区域。FF3R先求同一体素的平均语义向量$\bar{\bm{f}}_s$，以余弦距离$d_i^{\mathrm{sem}}$度量每项偏离，再把几何置信$C_i$与语义一致性共同用于聚合：
\begin{equation}
 \begin{aligned}
 w_{i\to s}=
 \frac{\exp(C_i-\lambda d_i^{\mathrm{sem}})}
 {\sum_{j:\,v(j)=s}\exp(C_j-\lambda d_j^{\mathrm{sem}})},\\
 \bar{\bm{a}}_s=\sum_{i:\,v(i)=s}w_{i\to s}\bm{a}_i\text{。}
 \end{aligned}
 \label{eq:vision-ff3r-voxel}
\end{equation}
$v(i)$是体素编号，$\bm{a}_i$可以是位置、颜色或语义属性，$\lambda$控制语义距离权重。权重和为1，所以它是体素内的加权汇总，不是把全部属性相加。若教学中两项置信均为1、语义距离分别为0与1、$\lambda=1$，权重为$e/(e+1)\approx0.731$与0.269；偏离项贡献下降但仍未消失。跨物体体素过大或平均语义本身错误，汇总仍会污染边界；平均语义向量为零时还需余弦的数值退化处理。

完整目标包含RGB像素与感知误差、特征余弦误差、双向几何特征误差以及深度和相机蒸馏项。按采样图集计算这些可微残差，更新可训练部分，训练完成后保留融合和解码参数；新房间使用前馈预测获得几何、外观与语义表示。查询“椅子”时，把文字编码为相同语义空间向量，与三维或投影后的特征比较，再按规定阈值和区域规则恢复结果。仅颜色渲染自然，不能证明椅子语义或尺度正确。

同一椅面的两次观察应落到共同三维位置并具有相容语义，未见背面则仍只有先验预测。检查可见表面的深度、跨视角语义归属、遮挡筛选和体素边界，才能判断查询结果是否真的指向同一对象。作者项目页标注CVPR 2026 Findings，这一身份不同于把预印本泛称为CVPR主会论文。新视角外观在\ref{sec:vision-generation-novel-view}节展开，未观测内容的创作归\ref{chap:vision-multimodal-generation}章。

\section{图像匹配与检索}
\label{sec:vision-recognition-retrieval}

识别已经把图像变成类别、范围和测量。检索还需要依据查询，从图库找出相关图像并排序。紧凑描述适合筛选候选，局部对应保留细节，几何约束再检查对应是否能共同解释视角变化。

\subsection{查询与相关性的确定}
\label{sec:vision-recognition-relevance}

“找类似的红色汽车”和“找同一辆车的其他照片”使用不同正例。前者容许不同实例，只要外观或语义相近；后者需要车身细节、车牌等区分身份的证据，并可能允许颜色因光照变化。若没有预先规定相关性，高相似度排序甚至可能朝错误目标优化。

查询可以是整图或裁剪区域。整幅照片包含建筑、树木和天空，局部查询只指定一块招牌；两者对应的容许背景、遮挡和视角变化不同。图库中的近重复图也可能使任务过于容易。评价前应固定查询范围、正例定义、忽略项、图库版本和干扰图规模，文字查询则由\ref{chap:vision-multimodal-understanding}章继续连接。

\subsection{候选图像的形成}
\label{sec:vision-recognition-candidates}

在大图库中逐一做精细匹配成本很高。全局向量先快速保留较少候选，局部特征再进入后续核验。近似索引可以加速向量搜索，但它引入的召回损失与表示本身的区分能力应分开检查。

\subsubsection{全局特征聚合与检索学习：GeM}
\label{sec:vision-recognition-gem}

空间特征需要汇总为固定维向量。平均池化平等计入各位置，最大池化只保留最强响应。\term{广义均值池化}（generalized-mean pooling, GeM）在两者之间调节：对第$k$通道的非负特征$x_{p,k}$，计算
\begin{equation}
 g_k=\left(\frac1n\sum_{p=1}^{n}x_{p,k}^{\,q_k}\right)^{1/q_k},
 \qquad q_k>0\text{。}
 \label{eq:vision-gem}
\end{equation}
指数$q_k=1$为平均，趋于无穷时趋近最大值。实现通常在幂运算前给特征加下界，避免零点数值问题；指数可以固定，也可与编码器参数一同学习\scite{vision-1711-02512}

对响应$(1,1,1,4)$，平均为1.75，$q=3$时为$(67/4)^{1/3}\approx2.559$，最大值为4。提高指数让局部强响应更影响全图描述，但并不保证该响应一定属于目标。聚合后将$\bm{g}$归一化为单位向量，以欧氏距离或内积排序；单位向量满足$\lVert\bm{u}-\bm{v}\rVert_2^2=2-2\bm{u}^{\top}\bm{v}$，两种排序因此等价。

检索训练需要知道哪些图拍到了同一实例。原工作利用三维重建与相机几何选择具有重叠内容的正例和困难负例，故“不需人工检索标注”并不等于没有监督依据。对比目标使正图对接近，对距离小于间隔$m$的负图对施加惩罚；训练后保留编码器与池化参数，为图库和查询使用相同的特征处理。

一条训练图对$(\bm X,\bm X')$的同实例标签为$y\in\{0,1\}$，两图经过共享编码、GeM和单位归一化后，距离为$d=\lVert\bar{\bm g}-\bar{\bm g}'\rVert_2$。对比目标的一种标准写法为
\begin{equation}
 \ell=\tfrac12 y d^2+
 \tfrac12(1-y)[\max(0,m-d)]^2\text{。}
 \label{eq:vision-gem-contrastive-process}
\end{equation}
正对的距离导数为$d$，下降更新倾向于缩短距离；负对在$d<m$时导数为$d-m<0$，倾向于增大距离，超过间隔则不给这项梯度。教学中$m=1$，正对距离0.4的目标为0.08，负对距离0.4的目标为0.18，负对距离1.2为零。平方距离是训练约束，不是与实际世界中建筑距离同量纲的测量。

原工作由重建得到的共同三维轨迹及相机共视挑选正对，从不同实例中挑选当前描述相近的困难负对。以已有分类编码器初始化，训练共同更新视觉特征与可学习池化指数，重新挖掘困难负例时使用当前参数，而测试图库不能参与形成监督。使用时为每图保存归一化描述与图号，对查询同样处理，按内积取前若干候选。指数、可能的描述变换和归一化须一致；只编码图库而忘记查询的同样变换，会改变距离定义。裁剪查询需保留裁剪来源，后续局部验证仍针对该查询范围。

同一建筑的两张照片可能因视角差异较远，另一栋同样玻璃幕墙的建筑却可能很近。全局描述提供候选，不能证明局部细节一致。已有DINOv3向量可以作为不同参数来源的筛选基线，但比较应固定输入、归一化、图库与训练信息。

\Needspace{55mm}
\subsubsection{局部关键点与描述：SuperPoint}
\label{sec:vision-recognition-superpoint}

若希望保留窗角或招牌角等可重复的位置，可以为每幅图输出稀疏点坐标和描述向量。SuperPoint使用共享编码器连接兴趣点分支与描述分支；兴趣点没有“左腕”那样的对象部位语义，而是适合跨图重复辨认的位置\scite{vision-1712-07629}

原始检测头在较粗网格上预测每个$8\times8$区域内的64个位置及一个“没有点”的通道。经过通道归一化、去掉无点通道和空间重排，得到原图分辨率的点分数，再阈值筛选并抑制过近峰值。描述头产生较稠密的低分辨率特征，在保留点位置插值采样，并将描述归一化。最终保存$(x,y,\text{得分},\bm{d})$，而不仅是一张响应图。

训练先以具有已知角点的合成几何图形学习检测，再对真实图像施加多次已知单应变换，将检测结果逆变换并汇总，形成更稳定的伪标签。已知变换也提供两图描述应对应的位置，用于联合训练点与描述。单应适应提供的是这些变换下的监督，不能覆盖任意三维视差。

对一个粗网格单元，64个点位置与无点通道组成65类，标签按变换后的可见兴趣点生成，交叉熵使点分数回到正确子位置。描述监督由已知单应把一个位置映到另一图，保留落在有效范围的对应；对应描述需内积较高，其他位置按间隔降低相似度。多次单应适应是产生真实图伪标签的步骤，联合训练再用这些伪标签及变换约束更新共享编码与两个分支。它与使用阶段在一对新照片上找对应不同，没有从新图真实相机取得标签。

得到点分数图后，按阈值与局部峰值抑制保留点，并可按分数限制最多点数；描述在这些点上插值并单位归一化。若输入缩放为原图的一半，网络点$(100,60)$需还原为原图$(200,120)$，描述向量本身则不乘二。保存点坐标所在网格和图号，后续匹配器才能用同一坐标系做几何验证；没有点时应返回空特征集合，不虚造一对对应。

窗格重复时，检测器可以在许多窗角上都稳定响应，但描述仍可能相似。因此“点能重复检测”与“点能匹配到正确窗格”是两个要求。尺度、旋转、遮挡及保留点数都会影响后续候选，比较匹配器时需固定这些输入条件。

\subsection{精细匹配与结果排序}
\label{sec:vision-recognition-verification}

候选图像中可能混入外观相似的其他实例。局部描述取得候选点对，上下文比较帮助消歧，几何一致性再排除不能共同成立的对应。几何通过也仍需回到查询定义判断相关性。

\Needspace{55mm}
\subsubsection{上下文局部配对与自适应计算：LightGlue}
\label{sec:vision-recognition-lightglue}

独立最近邻匹配只比较两个描述，无法利用邻近窗格怎样共同移动的信息。LightGlue接收已有关键点和描述，交替通过图内自注意力与跨图信息交换更新表示，再预测匹配关系；它不负责从像素检测关键点\scite{vision-2306-13643}

设两图分别有$n,m$个点，将更新后的描述投影为$d$维，再以$1/\sqrt d$缩放两两内积，得到$S_{i,j}$。两侧点可找到对应的分数为$\mu_i,\nu_j\in[0,1]$，匹配得分结合双向归一化与可匹配性：
\begin{equation}
 P_{i,j}=\mu_i\nu_j
 \frac{\exp S_{i,j}}{\sum_{j'=1}^{m}\exp S_{i,j'}}
 \frac{\exp S_{i,j}}{\sum_{i'=1}^{n}\exp S_{i',j}}\text{。}
 \label{eq:vision-lightglue-score}
\end{equation}
第一项归一化要求$i$在对面候选中偏好$j$，第二项要求$j$在本侧候选中偏好$i$。保留互为最高分且超过阈值的点对，有助于避免多个点争抢同一对应。训练使用真实对应监督匹配得分，对没有对应的点监督可匹配性。

一条匹配训练记录含两图的点坐标、描述和真实点对，以及哪些点在另一图不可见。点对可以由有深度、相机的图像重投影获得，也可以由已知图像变换取得；遮挡和越界必须排除，否则“找不到对应”的点会被错误配对。输入先按图像尺寸规范化关键点坐标，为点加入位置表示，再更新描述。若真实匹配集合为$\mathcal M$，两侧无对应点集为$U,V$，一个说明职责的负对数目标为
\begin{equation}
 \begin{aligned}
 L_{\mathrm{match}}={}&-\frac1{|\mathcal M|}
 \sum_{(i,j)\in\mathcal M}\log P_{i,j}\\
 &-\frac1{|U|}\sum_{i\in U}\log(1-\mu_i)
 -\frac1{|V|}\sum_{j\in V}\log(1-\nu_j)\text{。}
 \end{aligned}
 \label{eq:vision-lightglue-training-process}
\end{equation}
空集合对应项应跳过或用规定的非零归一化，不直接除零。匹配项同时监督双方向选择与可匹配性，无对应项监督放弃错误配对；它不是把全部$n m$点对都当同等负例。匹配器训练更新位置编码、信息交换和评分分支，关键点描述来源需固定，不能把为一种特征训练的权重无条件用于另一种描述。

若$\mu_i=0.9,\nu_j=0.8$，两个方向的归一化分数为0.7和0.6，则$P_{i,j}=0.3024$。只有单方向相似度高，或其中一点已被遮挡时，最终得分都可能降低。这是模型的配对评分，不是对几何正确性的独立证明。

不同图像对的难度差异很大。LightGlue在各层预测当前配对是否已稳定，足够比例的点达到置信要求时提前退出；对于已经高置信、但可匹配分数低的点，停止后续计算。尚不能判断的点仍应保留，不能把“不确定”直接当成“无对应”。这些门控调节处理层数和保留点数，比较速度时需固定特征、图像尺寸、硬件与难度分布。

层间置信分支以当前结果是否同最终层的决定相容为训练依据，并限制其梯度路径，避免仅为早退判断任意改动描述。使用时每层完成更新后检查达到置信要求的点比例，未达到则继续；被裁点保存原输入索引，最终对应需映射回完整点表。对保留点计算$P$，取行最大与列最大互相指向且得分过阈值的点对，再输出两幅原图的坐标及匹配得分。若任一图无点，应返回空对应；不能为了让后续几何估计运行而强行填点。置信早退和低可匹配裁减分别调节层数与点数，不改变后续还要做几何核验的职责。

\subsubsection{鲁棒几何估计与候选重排：RANSAC路线}
\label{sec:vision-recognition-ransac}

重复纹理可能产生看似可信但错误的点对，需要检查一组对应能否共同解释视角变化。平面上的点在两图间可由\term{单应变换}（homography）联系，齐次坐标满足$\tilde{\bm{p}}'\sim\bm{H}\tilde{\bm{p}}$。一般三维场景的两视图则遵守\term{极线约束}（epipolar constraint），即$\tilde{\bm{p}}'^\top\bm{F}\tilde{\bm{p}}=0$，其中$\bm{F}$为基础矩阵。后者约束另一图中的一条线，而非单个位置。

\term{随机采样一致性}（random sample consensus, RANSAC）通过反复取少量对应估计模型，寻找被许多点共同支持的解。以单应为例，四个处于非退化配置的点对可以产生候选$\bm{H}$；再把所有对应按重投影残差与阈值分为内点和外点，保存支持较好的模型，并利用内点重估。

几何模型的来源是当前点对，不需训练分类头。对平面对应，消去齐次比例可写$\tilde{\bm p}'\times\bm H\tilde{\bm p}=\bm0$，每对提供两条独立线性约束，四对约束八个尺度以外的自由度；共线、重复或近退化点会破坏估计。将输入坐标平移缩放到较合适的数值范围后求候选，再变回原图$\bm H$，才可用像素残差比较。对候选的对称转移误差可用
\[
 e_i^2=\lVert\bm p_i'-\Pi(\bm H\tilde{\bm p}_i)\rVert_2^2
       +\lVert\bm p_i-\Pi(\bm H^{-1}\tilde{\bm p}_i')\rVert_2^2\text{。}
\]
这里要求$\bm H$可逆且两次投影有限，阈值按像素或上述平方像素明确约定。一般三维两视图则不能用单应的点到点误差替代极线误差；$\bm F$只约束对应点靠近极线，估计时还需满足秩二条件。

若每次独立抽取$s$个对应，内点比例近似为$w$，全内点样本概率约为$w^s$。为了至少以概率$p$抽到一次全内点样本，次数满足
\begin{equation}
 k\geq\frac{\log(1-p)}{\log(1-w^s)}\text{。}
 \label{eq:vision-ransac-iterations}
\end{equation}
这是独立、近似固定内点率下的采样计算，还没有保证样本不退化，也没有保证估计模型正确。取$w=0.5,s=4,p=0.99$，需要至少72次；内点率降低时，所需次数迅速增加。

次数公式来自“$k$次都没抽到全内点”的概率$(1-w^s)^k$，要求它不超过$1-p$后取对数；分母为负数，除法时不等号方向随之改变。$w=0$时不存在有限次数保证，$w=1$时理想采样一步足够，实际仍需排除退化。实现从一个未接受模型状态开始，反复抽样、估计、计内点，若支持更多或在同支持下残差更小就替换当前最佳；到预算或按更新内点率估计的停止次数结束，再用全部内点重估。使用内点率提前停止依赖估计可信，错误高支持模型也可能过早结束。

Graph-Cut RANSAC在局部优化中进一步利用邻近对应的联系。为每个对应设置内外点标签，一元代价来自对当前模型的残差，邻接项倾向于让邻近对应具有一致标签；图割在规定的能量形式下求标记，再据选出的内点重估几何\scite{vision-1706-00984}
图割解决的是给定模型附近的一次局部标记问题，不是对所有模型、样本和标签组合取得全局最优。邻近错误匹配也可能具有一致外观，仍须检查模型假设。

在当前候选模型下，二值标签$\ell_i$表示该对应是内点还是外点，可把局部目标概括为$\sum_iD_i(\ell_i;e_i)+\lambda\sum_{(i,j)\in E}w_{i,j}\mathbf1\{\ell_i\ne\ell_j\}$。一元项依据残差惩罚不合适标记，非负邻接项表达邻近对应倾向于共同支持一个模型；在满足二值可割能量条件时，图割求这一次标记最优。重新用内点估计几何后，一元项改变，仍需继续局部比较。检索重排可以依据内点数量、查询范围覆盖和残差，但应固定排序规则；一小块重复窗格的高内点数不能替代整幅查询的相关性。

建筑正面近似平面时，单应适用；有明显前后视差时，同一个单应不能同时对齐所有表面，应考虑两视图几何或分区结构。纯旋转、近共线点和重叠太少也会改变估计条件。图\ref{fig:vision-retrieval-evidence}将全局筛选、局部配对与几何核验分别列出，避免把一个阶段的高分误作整条证据链完成。

\begin{figure}[htbp]
 \centering
 % 流程节点间箭头均连接左右中心；下方为各阶段尚未排除的错误。
\begin{tikzpicture}[font=\figurefont\small,
 block/.style={draw=BookTeal,fill=BookTeal!8,minimum width=28mm,
   minimum height=12mm,align=center,inner sep=1.5mm},
 flow/.style={-{Stealth[length=2mm]},draw=BookInk,line width=0.9pt},
 note/.style={align=center,text width=28mm,font=\figurefont\footnotesize,
   text=BookMuted}]
 \node[block] (global) at (0,0) {全局描述\\筛选候选};
 \node[block] (local) at (3.7,0) {局部对应\\匹配细节};
 \node[block] (geo) at (7.4,0) {几何支持\\核验重排};
 \draw[flow] (global.east)--(local.west);
 \draw[flow] (local.east)--(geo.west);
 \node[note] at (0,-1.6) {可能只同类\\并非同一实例};
 \node[note] at (3.7,-1.6) {重复窗格\\可能相互错配};
 \node[note] at (7.4,-1.6) {几何匹配仍需\\满足查询定义};
 \foreach \x in {0,3.7,7.4}
   \draw[draw=BookRule,line width=0.7pt] (\x,-0.65)--(\x,-1);
\end{tikzpicture}

 \caption{图像检索中逐步保留更具体的证据。上方沿主流程减少候选，下方说明每步仍可能留下的错误。几何一致可以帮助重排，却仍要检查所得实例是否满足查询的相关性定义。}
 \label{fig:vision-retrieval-evidence}
\end{figure}

Revisited Oxford和Paris修订了实例检索标注与难度协议，并引入大规模干扰图，适合比较候选筛选和几何重排分别改变了哪些结果\scite{vision-1803-11285}
不同协议对困难正例与忽略项的处理不一样，应固定协议后报告排序。更多内点可以支持同一建筑的判断，却不能直接证明其中招牌文字相同，或建筑具有相同的实际尺度。

\section{结果检验}
\label{sec:vision-recognition-evaluation}

组合系统需要让类别、框、掩码、属性和测量指向同一对象。若车辆检测正确，属性却来自旁车，逐模块的总体准确率可能仍然较高。街景应逐实例保留标识并核验空间对应；工业零件应同时检查异常图、整图得分与规定缺陷，避免正常材质变化被业务阈值放大。

检测评价先按类别与IoU阈值对预测和真实对象作一对一匹配，未匹配预测为假阳性，未匹配目标为漏检。按得分变化形成精确率、召回率曲线，再计算平均精度；不同数据集的插值、IoU阈值、尺寸分组与每图检测数量上限必须说明，不能把任意一个“AP”视为同一指标。

区域任务也有不同分母。语义分割按类别比较像素IoU，常再对类别平均；实例分割要先匹配对象。全景质量把匹配区域的IoU与对象错误合在一起，标准任务对同类、IoU超过0.5的区域建立匹配集合$T$，并令$F_{\mathrm{p}},F_{\mathrm{n}}$为多报、漏报区域集合：
\begin{equation}
 \mathrm{PQ}=
 \frac{\sum_{(p,g)\in T}\operatorname{IoU}(p,g)}
 {|T|+\frac12|F_{\mathrm{p}}|+\frac12|F_{\mathrm{n}}|}\text{。}
 \label{eq:vision-pq}
\end{equation}
它同时受到范围与实例错误影响\scite{vision-1801-00868}
例如匹配区域IoU为0.8和0.6，另有一个漏报且无多报，则PQ为$1.4/2.5=0.56$；只报告匹配区域平均IoU会得到0.7，掩盖漏掉的区域。空类别与跨类别平均按数据集协议处理。

检索则按相关性检查排名。若有$R$个相关项，排名$k$处的相关指示为$r_k$，截至该位置的精确率为$P(k)$，全排名上的一种非插值平均精度为
\begin{equation}
 \mathrm{AP}=\frac1R\sum_k P(k)r_k\text{。}
 \label{eq:vision-retrieval-ap}
\end{equation}
三个正例出现在第1、3、5位时，AP为$(1+2/3+3/5)/3\approx0.756$。若只返回前四项，计算截断指标时仍需明示分母和遗漏规则。召回率关注是否找到，AP还关注出现得早不早；无相关项的查询应按协议排除或单独评价，不能直接除以零。

未知发现需在隐藏新类别名称的条件下计未知召回，同时记录未知被误报成已知的数量、背景假阳性，以及后续加类前后旧类的AP。教学中有5个未知物体，正确发现3个、误当已知1个、完全漏掉1个，未知召回为$3/5=0.6$；这不能由已知货车AP代替。文字定位和概念分割还需固定候选短语、正负图像示例及提示成本，使用真实框提示取得的优势不能直接归于无提示系统。

姿态比较固定关节集合、单位和坐标，以平均关节位置误差（mean per-joint position error, MPJPE）计算$(1/j)\sum_k\lVert\hat{\bm J}_k-\bm J_k\rVert_2$。根对齐只去平移，相似变换对齐还去旋转和尺度，二者回答的问题不同。人体表面另比对应顶点误差、轮廓重投影及手足局部，不能由少量关节误差证明全部网格准确。单目深度需区分相对排序、边界和米制误差；允许对齐的结果不能与未对齐度量结果混为一表。

离线场景比较固定视图集合、分辨率、首相机参照和拟合步数，分别报告相机、点图、重投影、完整度及总成本。流式场景在每个时刻只允许已到达图，记录缓存预算、局部误差与长程漂移，回访时检查同一位置是否重合；一次全集对齐后的轨迹误差会隐藏部分累计误差。三维语义需先在有效几何下找到同一表面，再比较各视角的查询归属。RGB渲染、语言相似度和几何一致各有误差来源，三者之一正确不能替其他两项通过。

\begin{table}[htbp]
 \centering
 \small
 \begin{tabularx}{\linewidth}{@{}lXX@{}}
 \toprule
 输出 & 需要固定的条件 & 不能由它单独证明的内容\\
 \midrule
 类别、属性 & 类别表、未知项、对象范围 & 对象存在及属性归属都正确\\
 框、掩码 & 匹配规则、提示、可见范围 & 未观测部分及三维形状正确\\
 数量、姿态 & 计入规则、关节与对齐方式 & 实例清单、身份及原始朝向正确\\
 深度、点图 & 尺度、坐标、有效区域 & 距离带实际单位且适合测量\\
 检索排序 & 正例、图库、忽略项及难度 & 对象属性、文字和距离正确\\
 \bottomrule
 \end{tabularx}
 \caption{不同视觉成果所需的评价条件。右列列出仍需另行核验的推论，不表示这些任务必然无法提供相应能力。}
 \label{tab:vision-recognition-protocols}
\end{table}

\section{本章小结}
\label{sec:vision-recognition-summary}

从街景得到可用结果，需要把预测接回明确的对象和观测。类别与属性规定语义，框和掩码规定范围，数量与姿态规定测量对象，相机与尺度规定三维解释，检索相关性则决定哪些对应值得保留。共享表示可以支持这些输出，监督、恢复步骤和核验关系仍各自不可省略。

本章的方法提供了两类互补依据。学习表示利用样本规律形成候选判断，空间、几何与集合约束检查这些判断能否共同成立。候选高分不等于目标存在，区域正确不等于身份正确，相对深度不等于米制距离，几何支持也不等于任意业务相关性。只有把这些条件保留到最终输出，才能知道一项结果已经回答了什么，以及仍需要哪些观测。

开放词汇是在已给概念下寻找，未知发现是在未命名时保留对象，两者不能以同一查询准确率替代。二维关节、三维骨架和人体表面逐步增加几何约束，却始终受尺度与遮挡限制。离线全集压缩、因果滚动缓存又分别保存不同历史：固定参数容量不免除拟合和图像存储，有限缓存不保证长期回访没有漂移。三维语义只有与位置和可见表面共同核验，才构成可用于后续测量、渲染或问答的结果。

\Needspace{95mm}
这些尚未满足的条件也给出了进一步研究的问题。开放概念增加目标歧义，长尾类别需要更细的覆盖分析，精细边界要求保存小尺度证据，遮挡和单目投影又限制可恢复的空间量。LVIS把大词表和少见实例纳入分割评价，MVTec AD把正常训练与测试缺陷定位分开，为这些问题提供了不同的观察条件\scite{vision-1908-03195,vision-mvtec2019}
医学和遥感等专业用途还要固定采集设备、地理区域或患者划分，核验外观先验能否迁移。研究改进应明确补充了哪种证据、减少了哪类错误，结论仍限于相应任务与划分实际检验的范围。

遮挡人体的提示纠错还可研究如何表达剩余深度歧义，而非只输出一个看似合理的网格；长期地图可研究有限记忆、显式回访校正与尺度一致如何配合；几何语义融合可研究未见表面的预测如何同已观察事实分开保存。这些问题需要固定提示、视图到达、缓存及标注条件，分别检验关节、漂移和跨视角归属，不能用单一视觉分数宣称全部问题已解决。
